\documentclass[11pt]{article}

% ---------------------------------------------------------------------------
% Portable preamble: compiles with pdfLaTeX or XeLaTeX, no project macros.
% Mirrors doubletree/inst/paper/theory.tex's preamble style, since this
% paper's core content is being extracted from that document's Section 6
% (the two-stage estimator's topology-recovery / off-grid-refinement
% machinery, split out 2026-09-01 because it has no causal-inference content
% and stands as an independent contribution -- see
% quality_reports/plans/2026-09-01_two-paper-split.md at the repo root).
% ---------------------------------------------------------------------------
\usepackage{amsmath,amsthm,amssymb}
\usepackage[margin=1in]{geometry}
\usepackage[colorlinks=true,linkcolor=blue,urlcolor=blue,citecolor=blue]{hyperref}
\usepackage{mathtools}
\usepackage{enumitem}
\usepackage{booktabs}
\usepackage{longtable}
\usepackage[round,authoryear]{natbib}

\theoremstyle{plain}
\newtheorem{theorem}{Theorem}[section]
\newtheorem{lemma}[theorem]{Lemma}
\newtheorem{proposition}[theorem]{Proposition}
\newtheorem{corollary}[theorem]{Corollary}
\theoremstyle{definition}
\newtheorem{assumption}{Assumption}
\newtheorem{definition}[theorem]{Definition}
\theoremstyle{remark}
\newtheorem{remark}[theorem]{Remark}
\newtheorem{example}[theorem]{Example}

% NOTE: when porting content from theory.tex, also port the specific macro
% definitions that section actually uses (E, PP, argmin/argmax, cA_r,
% \topol, \topolcoll, \gbar, \cTcont, etc.) rather than pulling in
% common-defs.tex's full causal-nuisance vocabulary (\ehat, \thetahat, ...),
% which this paper does not need. A copy of common-defs.tex is available in
% this directory if a later section needs its general-purpose macros.
%
% The block below is exactly that: the macros actually used by the ported
% Sections~\ref{sec:setup}--\ref{sec:stageb}, transcribed from theory.tex's
% own preamble. Deliberately omitted, because this paper has no
% causal-nuisance objects: \ez, \mz, \wz, \ehat, \mhat, \what, \that,
% \thetahat, \thetatil, \pihat, \Vhat, \Dw, \Dm, \Lw, \Pn.
\newcommand{\E}{\mathbb{E}}
\newcommand{\PP}{\mathbb{P}}
\newcommand{\cX}{\mathcal{X}}
\newcommand{\cA}{\mathcal{A}}
\newcommand{\cB}{\mathcal{B}}
\newcommand{\cT}{\mathcal{T}}
\newcommand{\cS}{\mathcal{S}}
\newcommand{\cC}{\mathcal{C}}
\newcommand{\cG}{\mathcal{G}}
\DeclareMathOperator*{\argmin}{arg\,min}
\DeclarePairedDelimiter{\abs}{\lvert}{\rvert}
\DeclarePairedDelimiter{\card}{\lvert}{\rvert}

% --- Two-stage additions, matching theory.tex's own \S6 macro block ---
\newcommand{\topol}{\operatorname{top}}
\newcommand{\cTcont}{\cT^{\mathrm{cont}}_{\bar L}}
\newcommand{\gbar}{\bar\gamma}
\newcommand{\topolcoll}{\topol_{\mathrm{coll}}}

% Assumption numbering convention carried over from theory.tex: assumptions
% are numbered globally (A1, A2, ...), independently of the section-scoped
% theorem/lemma/proposition counter.
\renewcommand{\theassumption}{A\arabic{assumption}}

\title{Recovering Continuum-Threshold Optimal Trees:\\[2pt]
  \large A Two-Stage Algorithm with Superconsistent Threshold Refinement}
\author{[author list]}
\date{\today}

\begin{document}
\maketitle

\begin{abstract}
[TODO: abstract. Working content, per the 2026-09-01 split-planning
discussion: existing globally-optimal decision-tree solvers search a
\emph{pre-specified, fixed} grid of candidate cutpoints -- a genuine
threshold almost never lands exactly on one. We give a two-stage procedure
that removes this restriction without re-entering the combinatorial solver.
Stage A fits an ordinary grid-restricted optimal tree, at a coarse
resolution and an inflated leaf/depth budget, to recover the tree's
\emph{topology} (which coordinate splits where) rather than its exact
thresholds; a collapse map removes the ``transition leaf'' artifacts this
creates. Stage B then refines each threshold by an exact, grid-free
one-dimensional scan over observed covariate values, node by node,
root-first. We show topology is recovered with probability tending to one
on a refining grid, and that every refined threshold is simultaneously
superconsistent, at rate $O_p(n^{-1})$. The algorithm and its
guarantees are implemented and tested in the \texttt{optimaltrees} package.]
\end{abstract}

\tableofcontents

% ===========================================================================
\section{Introduction}\label{sec:intro}

[TODO -- outline, not drafted. Audience target revised 2026-09-01 to
JASA Theory \& Methods (was JMLR; see
\texttt{quality\_reports/plans/2026-09-01\_two-paper-split.md} \S7) --
lead with the statistical framing, not the ML motivation, per that
reconsideration. Candidate structure:
\begin{itemize}
\item Motivate (stats-first): decision-tree fitting is, at each split, a
  threshold/change-point estimation problem. Classical change-point theory
  (\citealt{chan1993}, \citealt{yaoau1989}, \citealt{baiperron1998},
  \citealt{boysen2009}, \citealt{seijosen2011}) and the split-point
  confidence-set literature (\citealt{banerjee2007}) study this for a
  single threshold, or for a tree used as a working (possibly
  misspecified) model. This paper asks the analogous question for a
  genuinely tree-structured target with multiple, jointly estimated
  thresholds arranged in a partition topology that must itself be
  recovered before any individual threshold's rate is even meaningful.
\item Secondary/practical motivation: existing globally-optimal
  decision-tree solvers (GOSDT, OSRT, this package's own branch-and-bound)
  search a \emph{pre-specified} grid of candidate cutpoints -- a genuine
  threshold generically does not land on a grid point. The two-stage
  algorithm below removes this restriction; the statistical question above
  is what makes the removal rigorous rather than heuristic.
\item Position against \citet{xu2026}: their oracle inequalities and
  minimax rates for ERM decision trees concern
  \emph{approximation}/prediction-risk theory over a nonparametric function
  class. This paper's question is different in kind -- whether the
  \emph{partition itself} is exactly recovered, and at what rate -- and is
  not addressed by that or related oracle-inequality results.
\item State the two contributions: (i) the two-stage algorithm; (ii) the
  topology-recovery + simultaneous-superconsistency guarantee.
\item Related work, expanded per the 2026-09-01 lit review
  (\texttt{quality\_reports/2026-09-01\_paper-a-lit-review.md}):
  \begin{itemize}
  \item Classical change-point/threshold estimation as above, plus
    \citealt{seootsu2018} (the argmax theorem for discontinuous criteria)
    and \citealt{song2016} (change-point asymptotics under
    misspecification) -- contrast: that literature gets cube-root,
    non-standard rates because the true regression function is smooth and
    the tree is a working (misspecified) model; here the true function is
    genuinely tree-structured, which is why superconsistency ($n^{-1}$,
    not $n^{-1/3}$) is available.
  \item ArborEnum (Heile, McTavish, Seltzer \& Rudin, 2026): the closest
    algorithmic cousin -- exact enumeration of Rashomon sets over
    continuous thresholds without prior binarization. Convergent
    motivation (binarization loses information), orthogonal contribution
    (combinatorial enumeration of a whole near-optimal set vs.\ this
    paper's statistical recovery-and-rate claim for one tree); neither
    subsumes the other.
  \item Xu (2026), \emph{On Stopping Rules and Spatial Adaptation for
    CART}: a companion piece (same lead author as \citealt{xu2026}) on
    greedy-CART pointwise rates -- adopt its global-vs.-pointwise rate
    vocabulary when scoping this paper's own (global) claim.
  \end{itemize}
\item Roadmap.
\end{itemize}
Bibkeys \texttt{song2016} and the ArborEnum/Xu-companion citations are
used above by name only (no \texttt{\textbackslash bibitem} yet), matching
how \texttt{xu2026} is already handled in this outline -- add all of them
together when this section is actually drafted.]

\section{Setup and the two-stage algorithm}\label{sec:setup}

\subsection{Covariate space, tree partitions, and the fitting problem}
\label{sec:setup-basics}

We observe $n$ independent copies of $(X,Z)$, where $X$ takes values in a
covariate space $\cX$ and $Z$ is a scalar response, and we wish to recover the
regression function
\[
  \gamma_0(x):=\E[Z\mid X=x],
\]
which is assumed throughout to be piecewise constant on an axis-aligned tree
partition of $\cX$. Both a regression problem (real-valued $Z$) and a
classification problem (binary $Z$, with $\gamma_0$ the conditional class
probability) are covered without change; the only structure used below is that
the fitting criterion is squared error, so that excess risk equals squared
$L_2$ distance to $\gamma_0$.

$\cX$ need not itself be finite; it may be any measurable space on which the
covariate law $P_X$ and the regression function $\gamma_0$ are well defined (a
finite product $\cX=\prod_{j\le d}\cX_j$ is the special case in which every
coordinate's cutpoint family below is its full value set). What is finite, and
fixed in $n$, is not $\cX$ but the \emph{candidate partition family},
generated by a pre-specified finite collection of cutpoints on each
coordinate --- chosen by the analyst before seeing the data and never updated
by it. These cutpoints partition $\cX$ into a finite family of \emph{grid
atoms} $\cA$, $\card\cA<\infty$ fixed in $n$, each atom a product of one
inter-cutpoint interval (or category group) per coordinate; write $P(a)$ for
the mass of $a\in\cA$. When $\cX$ is itself finite, atoms coincide with the
individual cells of $\cX$.

A \emph{tree partition} $\tau$ of $\cX$ is a partition induced by recursive
binary splits \emph{at the pre-specified cutpoints only}; its blocks are
\emph{leaves}, each leaf a union of atoms, $\card\tau$ is its leaf count, and
$\ell_\tau(x)$ is the leaf containing $x$. Restricting splits to the
pre-specified cutpoints is Assumption~\ref{ass:solver} below, not a
convention: without it the candidate class ceases to be fixed and finite. For
a leaf budget $\bar L$,
\[
  \cT_{\bar L}=\bigl\{\tau:\tau\text{ a tree partition of }\cX\text{ built
  from the pre-specified cutpoints},\ \card\tau\le\bar L\bigr\},
\]
a fixed finite set: any $\tau\in\cT_{\bar L}$ is a partition of the finite set
$\cA$ into at most $\bar L$ blocks realisable by cutpoint splits, so
$\card{\cT_{\bar L}}<\infty$ and does not depend on $n$. More generally, for a
finite cutpoint family generating atoms $\cB$ we write $\cT_\cB$ for the
unrestricted class of tree partitions built from those cutpoints, and
$\cT_{L}(\cB)$ for its restriction to at most $L$ leaves; thus
$\cT_{\bar L}=\cT_{\bar L}(\cA)$.

The \emph{topology} of a tree partition is its split structure with the
threshold \emph{values} discarded: formally, the rooted binary tree together
with the splitting coordinate attached to each internal node, and the
ancestor--descendant order of those coordinate tests. Two trees share a
topology when they differ only in where, along the designated coordinates,
their thresholds sit. Recovering the topology is therefore a strictly weaker
task than recovering the partition, and it is the task Stage~A below is asked
to perform.

For a strictly positive weight $\nu$ on $\cX$, let $\Pi^\nu_\tau$ denote the
$\nu$-weighted projection onto functions constant on the leaves of $\tau$;
we write $\Pi_\tau$ for $\Pi^{P_X}_\tau$. The measure $\nu$ is the
\emph{design measure} of the fitting problem: it is a fixed, strictly
positive measure on $\cX$, equivalent to $P_X$ in the sense of
Assumption~\ref{ass:design} below, and every mass, projection, and risk
below is taken with respect to it. Taking $\nu=P_X$ recovers the ordinary
unweighted problem; allowing a general $\nu$ costs nothing in the arguments
and covers weighted least squares, in which the analyst's target is the
$\nu$-weighted rather than the $P_X$-weighted projection of $\gamma_0$.

Given $\tau$, the \emph{full-sample leaf refit} sets, on each leaf $\ell$,
\begin{equation}\label{eq:refit}
  \widehat\gamma_\tau\restriction_\ell
  =\frac{\sum_{i:X_i\in\ell}Z_i}{\card{\{i:X_i\in\ell\}}} .
\end{equation}
Structure is selected by penalised empirical risk minimisation over the
\emph{feasible} candidates. Let
\begin{equation}\label{eq:feasible}
  \cT_n=\bigl\{\tau\in\cT_{\bar L}:\text{every leaf of }\tau\text{ carries
   at least }m_n\text{ observations}\bigr\},
\end{equation}
and for a penalty $\lambda_n>0$ set
\begin{equation}\label{eq:select}
  \widehat\tau\in\argmin_{\tau\in\cT_n}
  \bigl\{R_n(\tau)+\lambda_n\card\tau\bigr\},
\end{equation}
where $R_n$ is the empirical risk of the leaf-refit fit and $R$ its population
counterpart. Both risks are on the unnormalised scale, so that $R(\tau)$ is
the $\nu$-weighted risk $\E_\nu[\{\Pi^\nu_\tau\gamma_0-\gamma_0\}^2]$ rather
than any conditional average of it. Restricting the argmin to $\cT_n$ is what
reconciles selection with the minimum-leaf-mass requirement of
Assumption~\ref{ass:construct}(d): an unconstrained argmin over
$\cT_{\bar L}$ could return a partition violating it, so the constraint
belongs in the selection rule rather than being asserted of its output. The
condition $m_n/n\to0$ follows from $\bar Lm_n\lesssim n$ with $\bar L$ fixed.

\subsection{Standing assumptions}\label{sec:setup-assumptions}

\begin{assumption}[Sampling]\label{ass:iid}
$(X_1,Z_1),\dots,(X_n,Z_n)$ are independent and identically distributed copies
of $(X,Z)\sim P$.
\end{assumption}

\begin{assumption}[Design measure and boundedness]\label{ass:design}
The design measure $\nu$ is a fixed measure on $\cX$, not depending on $n$,
for which there exist constants $0<c\le C<\infty$ with
\[
  c\,dP_X\ \le\ d\nu\ \le\ C\,dP_X .
\]
Moreover $\gamma_0$ is bounded: $\sup_{x\in\cX}\abs{\gamma_0(x)}\le
C_\gamma<\infty$.
\end{assumption}

\begin{remark}[What replaces an identification condition here]
\label{rem:design-role}
Assumption~\ref{ass:design} is used for exactly two purposes below, and for
nothing else. The two-sided comparison of $\nu$ with $P_X$ ensures that every
grid atom of positive $P_X$-mass has positive $\nu$-mass, which is what makes
the population ANOVA identity of Lemma~\ref{lem:grid-optimal} sharp; and the
bound on $\gamma_0$ ensures that the profiled leaf means appearing in
Lemma~\ref{lem:descendant-stability} are bounded, uniformly in the candidate
threshold. Neither use requires $\nu$ to arise from any particular modelling
consideration. In an application in which the design measure is itself
estimated, $\nu$ should be read as the population weight; the arguments below
are stated at a fixed $\nu$.
\end{remark}

\begin{assumption}[Moments]\label{ass:moments}
$\E[Z^2\mid X=x]\le C_Z<\infty$ for every $x\in\cX$. Equivalently, in view of
Assumption~\ref{ass:design}, the conditional residual second moment is
uniformly bounded:
\[
  \sup_x\E\bigl[\{Z-\gamma_0(X)\}^2\mid X=x\bigr]<\infty .
\]
\end{assumption}

\begin{assumption}[Finite candidate partition family]\label{ass:finite}
The grid atoms $\cA$ generated by the pre-specified cutpoints are finite in
number, fixed in $n$, and each of positive mass: $\card\cA<\infty$ and
$P(a)>0$ for every $a\in\cA$. ($\cX$ itself is not required to be finite; see
\S\ref{sec:setup-basics}.)
\end{assumption}

\begin{assumption}[Estimator construction]\label{ass:construct}
The fit is constructed as follows.
\begin{enumerate}[label=(\alph*),nosep]
  \item \emph{Loss.} The fit is by squared error, so that excess risk equals
  squared $L_2(\nu)$ distance to $\gamma_0$.
  \item \emph{Leaf refit.} Leaf values are given by~\eqref{eq:refit}: the
  within-leaf response mean.
  \item \emph{Range restriction.} Reported fitted values are confined to a
  fixed compact interval containing the range of $\gamma_0$, whose existence is
  Assumption~\ref{ass:design}. Structure selection uses the
  \emph{unrestricted} leaf refit; the restriction is applied only to the
  reported fit.
  \item \emph{Leaf mass.} Selection is restricted to the feasible set
  \eqref{eq:feasible}: candidates whose every leaf carries at least $m_n$
  observations, with $\bar Lm_n\lesssim n$. This is a property of the
  selection rule, not a claim about its output.
\end{enumerate}
\end{assumption}

\begin{assumption}[Solver fidelity]\label{ass:solver}
The optimisation procedure computing~\eqref{eq:select} considers only tree
partitions whose splits lie at the pre-specified cutpoints of $\cX$: the
search space is $\cT_{\bar L}$, never a data-dependent or
continuously-refined set of thresholds. When $\cX$ is itself finite with every
value a pre-specified cutpoint, this holds automatically for any solver, since
every threshold on a finite value set already lies at a cutpoint; it is
recorded as a standing assumption because it is not automatic once $\cX$ is a
general space and a generic branch-and-bound solver would otherwise consider
splits at observed data values, which are neither fixed nor finite in $n$.
Relaxing precisely this restriction --- without paying for it in the
combinatorial search --- is the purpose of the two-stage procedure of
\S\ref{sec:setup-twostage}.
\end{assumption}

\begin{assumption}[Fixed leaf budget]\label{ass:budget}
$\bar L$ is a constant, fixed in $n$, with $\bar L\le\card\cA$; consequently
$\cT_{\bar L}$ is a fixed finite set and $\card{\cT_{\bar L}}$ does not depend
on $n$.
\end{assumption}

\begin{assumption}[Exact global optimisation]\label{ass:global}
$\widehat\tau$ attains the global minimum in~\eqref{eq:select} over the
feasible set $\cT_n$ of~\eqref{eq:feasible}, and $\lambda_n\to0$.
\end{assumption}

\begin{definition}[Sufficient class]\label{def:sufficient}
\[
  \cS=\bigl\{\tau\in\cT_{\bar L}:
  \Pi^\nu_\tau\gamma_0=\gamma_0\ \ P\text{-a.s.}\bigr\},
\]
equivalently the set of $\tau\in\cT_{\bar L}$ on whose leaves $\gamma_0$ is
$P$-a.s.\ constant. The characterisation is independent of the strictly
positive weight $\nu$.
\end{definition}

\begin{assumption}[Structural sparsity]\label{ass:sparsity}
$\cS\ne\emptyset$; equivalently, some tree partition built from the
pre-specified cutpoints, with at most $\bar L$ leaves, represents $\gamma_0$
exactly.
\end{assumption}

Assumption~\ref{ass:sparsity} is the grid-exact hypothesis this paper is
designed to dispense with: it requires every true threshold to coincide with a
pre-specified cutpoint. Assumption~\ref{ass:sparsity-jump}(S) below replaces
it by its continuum analogue.

\subsection{Continuum sparsity and the two-stage estimator}
\label{sec:setup-twostage}

The class-relativised form of Definition~\ref{def:sufficient} will be used
below. For a general tree class $\cC$, define
\[
  \cS(\cC)
  :=
  \bigl\{\tau\in\cC:
  \Pi^\nu_\tau\gamma_0=\gamma_0\ \ P\text{-a.s.}\bigr\}.
\]
Thus $\cS=\cS(\cT_{\bar L})$ denotes the grid-restricted class of
Definition~\ref{def:sufficient}, whereas
\[
  \cS^{\mathrm{cont}}
  :=
  \cS(\cTcont)
\]
denotes its continuum analogue, $\cTcont$ being the class of axis-aligned tree
partitions with at most $\bar L$ leaves whose thresholds are unrestricted real
numbers. Clause (S) of Assumption~\ref{ass:sparsity-jump} asserts
$\cS^{\mathrm{cont}}\ne\emptyset$. This is strictly weaker than the grid-exact
Assumption~\ref{ass:sparsity}, because a true continuum boundary need not
coincide with any cutpoint of a finite grid.

It is worth stating at the outset what this paper borrows and what it adds,
since the two are easily conflated. The localisation machinery is not new: the
rate-$n$ behaviour of a least-squares change point under random design is due
to \citet{chan1993} and \citet{seijosen2011}; the simultaneous estimation of
finitely many step heights alongside their locations is \citet{baiperron1998},
\citet{yaoau1989}, and \citet{boysen2009}; and the argmax theorem for
discontinuous, locally-defined criteria that converts a drifting local process
into a tight minimiser is \citet{seootsu2018}. The profiling-correction
argument of Remark~\ref{rem:profiling-correction} is likewise a standard
device, applied here rather than developed. Confidence sets for split points
have also been studied directly, by \citet{banerjee2007}, but under model
misspecification (a smooth true regression function approximated by a tree),
which yields non-standard, cube-root asymptotics; that is a different regime
from the one considered here, where the truth is exactly a tree, and it is
what makes the $n^{-1}$ threshold rate of
Proposition~\ref{prop:twostage-superconsistency} available at all.

What the tree structure adds is narrower, and amounts to two reductions and
one accounting device. First, Lemma~\ref{lem:prism-profile} shows that
refining one threshold inside a tree is not a two-group problem: the moving
slab must be decomposed into source--target leaf pairs, and the resulting
finite sum of marked processes --- not any single change-point process ---
is what carries the drift. Second, Lemma~\ref{lem:ancestor-localisation}
shows that superconsistency survives conditioning on \emph{estimated}
rather than true ancestors, which has no counterpart in a fixed-design
break model and is what makes the levelwise induction legitimate. Third,
Proposition~\ref{prop:transition-leaves} and the collapse map
$\topolcoll$ supply the budget accounting that lets a fixed-grid search
recover a continuum topology at all, by paying one extra leaf per
grid-misaligned boundary. Everything else below is the cited machinery
instantiated at those three points.

\begin{assumption}[Continuum tree sparsity and jump regularity]
\label{ass:sparsity-jump}
The regression function \(\gamma_0\) satisfies the following conditions.
\begin{enumerate}[label=(\Alph*),leftmargin=*]
  \item[(S)] \emph{Tree realisability.}
  The function \(\gamma_0\) is piecewise constant on a continuum
  partition realised by an axis-aligned tree with at most \(\bar L\) leaves
  and depth at most \(d\), where \(\bar L\) and \(d\) are fixed in \(n\).

  \item[(J)] \emph{Genuine jumps.}
  Every true split separates regions on which the relevant population values
  differ by an absolute amount between fixed constants
  \(0<\kappa_{\min}\le\kappa_{\max}<\infty\).

  \item[(Dens)] \emph{Local connected density.}
  Let \(R_N\) be the closed axis-aligned box associated with a true internal
  node \(N\), and let
  \[
    R_N^+
    :=
    \bigl\{x\in\cX:
      \operatorname{dist}_\infty(x,R_N)\le\varsigma/2
    \bigr\}.
  \]
  On \(R_N^+\), the covariate law under the measure \(\nu\) admits a density
  \(f_{\nu,N}\), with respect to Lebesgue measure or the appropriate product
  base measure, satisfying
  \[
    0<f_{\min}
    \le f_{\nu,N}(x)
    \le f_{\max}<\infty
    \qquad\text{for every }x\in R_N^+,
  \]
  where \(f_{\min}\) and \(f_{\max}\) are fixed uniformly over true nodes.
  (Every cross-sectional-measure argument below, in particular the
  definition of \(c_{\mathrm{box}}\) preceding Lemma~\ref{lem:topol-margin},
  is written for the Lebesgue case; a discrete or mixed splitting
  coordinate would in any case violate the \(\omega_{N,r}(c)=\Theta(\abs{c-c_0})\)
  prism scaling of Lemma~\ref{lem:prism-profile}(iii) on which the
  \(n^{-1}\) rate below depends, so this is a scoping remark, not an
  additional restriction.)
  The sets \(R_N^+\) are fixed compact rectangles whose coordinate diameters
  are bounded uniformly over true nodes by a fixed constant
  \(D_{\max}<\infty\).  In particular, the density lower bound implies
  that the support has no gaps in a neighbourhood of any true splitting
  face.  Thus, whenever \(x^-\) and \(x^+\) differ only in the splitting
  coordinate and lie sufficiently close to opposite sides of a true
  boundary, both points belong to the local support.  This condition rules
  out, by construction, supports consisting locally of disjoint rectangles.
  It is the local version of the customary assumption that covariates have a
  density bounded away from zero and infinity on a compact connected
  support.  The converse implication need not hold: connected support alone
  does not require its density to be bounded away from zero.

  \item[(Sep)] \emph{Separation and interiority.}
  Distinct true boundaries on the same coordinate are separated by at least
  \(\varsigma>0\), and every true boundary is at least \(\varsigma\) from the
  relevant support boundary.

  \item[(Node)] \emph{Positive node mass.}
  Every true node and leaf has \(\nu\)-mass at least \(q_{\min}>0\), fixed in
  \(n\).

  \item[(Local-Jump)] \emph{Jumps across every active face cell.}
  Fix an internal node \(N\), splitting coordinate \(q\) at \(c_{0,N}\). For
  every pair of leaves adjacent across this boundary, the corresponding
  population values \(\gamma^-_{N,r}\) and \(\gamma^+_{N,r}\) satisfy
  \(\kappa_{\min}\le\abs{\gamma^+_{N,r}-\gamma^-_{N,r}}\le\kappa_{\max}\).

  \item[(Joint)] \emph{Ancestor--descendant joint density.}
  If an ancestor split uses coordinate \(q_a\), a descendant split uses
  \(q_\ell\), and \(q_a\ne q_\ell\), the pair \((X_{q_a},X_{q_\ell})\) has a
  joint density bounded above by a fixed constant \(\bar f_2<\infty\) on a
  neighbourhood of \((c_{0,a},c_{0,\ell})\).
\end{enumerate}
\end{assumption}

\begin{remark}[Interpretation of the local-jump clause]
Condition (Local-Jump) is stronger than requiring merely that a split produce
some aggregate risk reduction.  It rules out cancellation caused by an
adjacent pair of leaves having the same population value.  Any design whose
true leaf values are pairwise distinct satisfies the condition by
construction, with \(\kappa_{\min}\) the smallest and \(\kappa_{\max}\) the
largest gap between values on leaves adjacent across a true face.
\end{remark}

For each resolution \(r\), let \(\cA_r\) be the full product grid generated by
\(r\) cutpoints per coordinate, up to fixed multiplicative constants: the
ratio of the largest to the smallest cutpoint spacing, in any coordinate and
at any resolution \(r\), is bounded by a fixed quasi-uniformity constant
\(\gamma_{\mathrm{grid}}\ge1\).  Its
mesh is \(O(r^{-1})\), tends to zero as \(r\to\infty\), and
\[
  N_r:=\card{\cA_r}\asymp r^d.
\]
Stage~A globally minimises the penalised empirical risk over trees built from
\(\cA_r\), subject to the inflated leaf budget
\[
  \bar L_A:=\min\{2\bar L-1,N_r\}
\]
and depth bound \(2d\) (the leaf budget is truncated at the grid's own
capacity \(N_r\) for small \(r\); the depth bound needs no matching
truncation, since a tree respecting \(\bar L_A\) already has depth at most
\(\bar L_A-1\le N_r-1\)).  These bounds accommodate the fact that a true
boundary lying strictly inside one grid atom may require two consecutive
grid splits, rather than one continuum split: the first split enters the
straddling atom and the second exits it.  This introduces at most one
additional leaf and one additional level for each true boundary.  Since a
tree with \(L\) leaves has \(L-1\) boundaries, at most \(2\bar L-1\) leaves
and depth \(2d\) are required.

Only the collapsed topology \(\topolcoll(\hat t^A)\) is retained from
Stage~A.  More precisely, suppose that within a node a split at cutpoint
\(u\) is followed, in the child containing the adjacent grid interval, by a
split on the same coordinate at the consecutive cutpoint \(v\).  If the
intermediate leaf is the transition atom between \(u\) and \(v\), the
collapse map replaces these two splits by the single labelled interval
\([u,v]\); all unmatched splits are retained as singleton intervals.  The
distance between a collapsed split and \(b\) is the distance from \(b\) to
its associated singleton or closed interval.  Thus a transition pair
straddling \(b\) has distance zero from \(b\).

In Stage~B, each recovered collapsed split is refined by an exact scan over
the observed values of its splitting coordinate.  The scan is restricted to
the bracket bounded by the nearest \emph{other} same-coordinate split in the
recovered topology, or by the support boundary if no such split exists.
The topology, descendants, and all other split coordinates are held fixed
during this scan.  Both stages are made precise in
Assumption~\ref{ass:solver-twostage} below.

If the minimum-leaf requirement of
Assumption~\ref{ass:construct}(d) is imposed in Stage~A, it applies with the
inflated budget \(\bar L_A\).  Transition leaves have population mass
\(\Theta(r^{-1})\); consequently their asymptotic feasibility requires the
additional condition
\[
  m_n=o(n/r_n).
\]
The existing condition \(\bar L_A m_n\lesssim n\) alone does not imply this
transition-leaf requirement.

\subsection{Solver fidelity for the two stages}\label{sec:setup-solver}

\begin{assumption}[Two-stage solver fidelity]\label{ass:solver-twostage}
The two-stage estimator satisfies the following conditions.
\begin{enumerate}[label=(\alph*),nosep]
  \item \emph{Stage-A optimisation.} Stage A computes an exact global minimiser
  of the penalised empirical risk over the grid $\cA_r$, using the inflated leaf
  and depth budgets. Every Stage-A split lies at a pre-specified cutpoint in
  $\cA_r$; thus the pre-specified-cutpoint requirement of
  Assumption~\ref{ass:solver} remains in force for Stage A itself.
  \item \emph{Stage-B optimisation.} Conditional on the topology returned by
  Stage A, Stage B refines thresholds one node at a time, in a fixed
  root-first order: at each node, it computes an exact global minimiser of
  the empirical risk over the observed covariate values lying in that
  node's own bracket (the interval bounded by the nearest other
  same-coordinate split in the recovered topology), holding every other
  threshold in the tree fixed --- already refined for every ancestor (the
  root-first order visits every ancestor before any descendant), still at
  its Stage-A grid value for every descendant and for any not-yet-visited
  node. Stage B is therefore a sequence of single-threshold exact scans, not
  a joint minimisation over the Cartesian product of all thresholds at
  once. Only Stage B may evaluate continuous or observed-value thresholds.
  \item \emph{Full-sample construction.} Stages A and B both use all $n$
  observations. There is no sample splitting or cross-fitting between them.
  \item \emph{Topology preservation.} Stage B refines threshold locations only:
  no split is added or deleted, no pair of splits is reordered, and no split is
  moved to a different covariate coordinate.
  \item \emph{Valid brackets.} For every recovered split, the bracket bounded by
  the nearest other split on the same coordinate on each side, as determined
  by the recovered topology, contains the corresponding true threshold with
  probability tending to one.
  \item \emph{Measurability.} Both stages use deterministic, measurable
  tie-breaking rules. Consequently the refined threshold vector $\hat c$ is a
  well-defined random vector.
  \item \emph{Refined leaf mass.} There exists a constant $\rho>0$, independent
  of $n$, such that, with probability tending to one, every leaf of the refined
  partition contains at least $\rho n$ observations. The same bound holds at
  the true partition $c_0$. Thus finite-dimensional leaf-refit control is
  available at the anchor and throughout the topology-preserving perturbation
  from $c_0$ to $\hat c$.
\end{enumerate}
\end{assumption}

\begin{remark}\label{rem:bracket-consequence}
Clause~(e) is a consequence of Proposition~\ref{prop:boundary-recovery}.
Indeed, on
\[
  E_n:=
  \Bigl\{
    \topolcoll(\hat t^A)\text{ is in bijection with the true boundaries of }
    \gamma_0\text{, each within }C/r_n\text{ of its match}
  \Bigr\},
  \qquad \PP(E_n)\to1,
\]
the recovered split corresponding to each true boundary lies within
$C/r_n$ of that boundary, while boundary separation and interiority from
Assumption~\ref{ass:sparsity-jump} prevent adjacent true boundaries from
entering the same bracket. No rate condition linking $r_n$ to $n$ is needed
beyond the conditions already required by
Proposition~\ref{prop:boundary-recovery}. Stage A supplies only the topology
and brackets; every threshold value in the final estimator is supplied by the
off-grid Stage-B scan.
\end{remark}


\section{Topology recovery on a refining grid}\label{sec:topology}

\noindent\emph{A note on reading Assumptions~\ref{ass:finite},
\ref{ass:budget}, and~\ref{ass:global} in this section.} Those assumptions
were stated in \S\ref{sec:setup} for the fixed grid $\cA$ and the fixed
budget $\bar L$. Every invocation below is at the resolution-$r$ grid
$\cA_r$ (finite for each fixed $r$, exactly as required) with the inflated,
still-fixed-for-each-$r$ budget $\bar L_A$; nothing below lets $r$ or the
budget vary within a single application of any of the three assumptions.
Where the text below writes, e.g., "Assumption~\ref{ass:finite}," it means
"Assumption~\ref{ass:finite}, read at $\cA_r$, for the $r$ fixed in that
application" throughout; this is restated at
Lemma~\ref{lem:grid-optimal} but not repeated at every subsequent use.

\begin{lemma}[The grid-optimal class is the projected-sufficient class]
\label{lem:grid-optimal}
Under Assumptions~\ref{ass:design} and~\ref{ass:finite} at $\cA_r$,
$\cS^{\cA_r}=\argmin_{\tau\in\cT_{\cA_r}}R(\tau)$, and
$\tau^{\mathrm{sat}}_r\in\cS^{\cA_r}$, so $\cS^{\cA_r}\ne\emptyset$
unconditionally, where $\tau^{\mathrm{sat}}_r$ is the partition of $\cA_r$
into its individual atoms,
$\gbar_r:=\Pi^\nu_{\tau^{\mathrm{sat}}_r}\gamma_0$
and
$\cS^{\cA_r}:=\bigl\{\tau\in\cT_{\cA_r}:\Pi^\nu_\tau\gbar_r=\gbar_r\
P\text{-a.s.}\bigr\}$.
\end{lemma}
\begin{proof}
The population ANOVA identity gives, for every $\tau\in\cT_{\cA_r}$,
$R(\tau)-R(\tau^{\mathrm{sat}}_r)
 =\sum_{\ell\in\tau}\sum_{a\subseteq\ell}\omega_a(\gbar_a-\gbar_\ell)^2\ge0$,
with $\omega_a$ the $\nu$-mass of atom $a$ and $\gbar_\ell$ the
$\omega$-weighted mean of $\gbar_a$ over $a\subseteq\ell$. Every atom has
$\omega_a>0$ by Assumptions~\ref{ass:finite} and~\ref{ass:design}, so the sum
vanishes exactly when $\gbar_r$ is leafwise constant. The membership of
$\tau^{\mathrm{sat}}_r$ is immediate.
\end{proof}

\paragraph{Proof audit.}
The argument uses only the population ANOVA identity, positivity of every
atom mass, and finiteness of the grid.  No stochastic rate or tail bound is
used.  It can fail if a purported grid atom has zero \(\nu\)-mass.

Fix \(r_0\), depending only on \(\varsigma\), the quasi-uniformity constant
\(\gamma_{\mathrm{grid}}\),
and the constant \(C\) selected below, such that, for every \(r\ge r_0\),
the mesh in every coordinate is smaller than \(\varsigma\) and
\(C/r<\varsigma/2\).  Hence no grid atom contains more than one true boundary.
This convention remains in force in the next three results.

A tree \(\tau\) is said to \emph{omit} a true boundary
\(b=(q,c_{0,N})\) if no singleton or interval in the collapsed split set
\(\topolcoll(\tau)\), on coordinate \(q\), lies within distance \(C/r\) of
\(c_{0,N}\).  Let \(D_{\max}\) be the coordinate-diameter bound of (Dens) and
set
\[
  c_{\mathrm{box}}:=\frac1{f_{\max}D_{\max}}>0,
\]
so that, for every true internal node \(N\) with splitting coordinate \(q\),
the \((d-1)\)-dimensional cross-section \(\Gamma_N\) of \(R_N\) transverse to
\(q\) satisfies \(\operatorname{Leb}_{d-1}(\Gamma_N)\ge c_{\mathrm{box}}
q_{\min}\).  Indeed (Node) gives \(\nu(R_N)\ge q_{\min}\) while (Dens) gives
\(\nu(R_N)\le f_{\max}\operatorname{Leb}_{d-1}(\Gamma_N)D_{\max}\).  Set
\[
  c_{\mathrm{topol}}:=\frac{c_{\mathrm{box}}}{8}.
\]
Choose the fixed constant \(C\) sufficiently large that
\begin{equation}\label{eq:C-margin-choice}
  c_{\mathrm{topol}}\kappa_{\min}^2 f_{\min}q_{\min}C
  >
  4(\bar L-1)\kappa_{\max}^2f_{\max}.
\end{equation}

\begin{lemma}[A refining-grid margin for omitted true boundaries]
\label{lem:topol-margin}
Under Assumptions~\ref{ass:design}, \ref{ass:budget}, and
\ref{ass:sparsity-jump}, there exist \(r_0<\infty\) and
\(c_\Delta>0\), independent of \(r\) and \(n\), such that, for every
\(r\ge r_0\) and every
\(\tau\in\cT_{\bar L_A}(\cA_r)\) that omits at least one true boundary,
\[
  S_r(\tau)
  :=
  R(\tau)
  -
  \min_{\sigma\in\cT_{\bar L_A}(\cA_r)}R(\sigma)
  \ge \frac{c_\Delta}{r}.
\]
One may take
\[
  c_\Delta
  :=
  c_{\mathrm{topol}}\kappa_{\min}^2f_{\min}q_{\min}C
  -
  (\bar L-1)\kappa_{\max}^2f_{\max}
  >0.
\]
\end{lemma}

\begin{proof}[Proof sketch]
Both halves of the bound are population-level; neither a concentration
inequality nor a union bound is used. For the lower bound, fix the omitted
boundary and consider a pair of points differing only in the splitting
coordinate and straddling it at distance $t\le C/2r$: no split on another
coordinate can separate them, and, precisely because the boundary is omitted,
no $q$-split of $\tau$ lies between them, so they share a leaf. Clause
(Local-Jump) then forces a squared error of at least $\kappa_{\min}^2/4$ at one
of the two points, and integrating over $t$ and over the active face cells ---
using the density floor of (Dens) together with the cross-sectional bound
$\operatorname{Leb}_{d-1}(\Gamma_N)\ge c_{\mathrm{box}}q_{\min}$ that (Node)
and (Dens) supply --- gives an approximation loss of order $C/r$. For the
upper bound, snapping every true boundary to a nearest cutpoint while
preserving the true ancestor order produces a feasible comparator in
$\cT_{\bar L_A}(\cA_r)$ that disagrees with $\gamma_0$ only in slabs of width
$O(r^{-1})$, costing at most $(\bar L-1)\kappa_{\max}^2f_{\max}/r$ by (J) and
the upper density bound. Subtracting the second bound from the first leaves
$c_\Delta/r$, and the choice~\eqref{eq:C-margin-choice} of $C$ is exactly what
makes $c_\Delta>0$ after the $O(r^{-1})$ discard loss at the edge of the face
patch is absorbed. The full proof is given in the Supplementary Material,
\S\ref{sec:supp-topology}.
\end{proof}

\begin{remark}[Population scope of the refining-grid margin]
Lemma~\ref{lem:topol-margin} is entirely population-level.  Its proof uses
neither a concentration inequality nor a union bound.  The
empirical-design comparison required for
Stage~A is separate and is established in
Proposition~\ref{prop:parsimony-grid}.
\end{remark}

\paragraph{Proof audit.}
The lower bound uses the density lower bound precisely to ensure that paired
points on both sides of a true face exist with uniformly positive local
mass.  The upper bound uses the density upper bound and the fixed number of
true boundaries.  Separation ensures that the paired patch contains no
second true boundary.  The conclusion can fail if the local support has a
gap, if a jump vanishes on an active face cell, or if node cross-sections
degenerate with \(n\).

\begin{proposition}[Parsimonious topology recovery on a refining grid]
\label{prop:parsimony-grid}
Suppose Assumptions~\ref{ass:iid}, \ref{ass:design},
\ref{ass:moments}, \ref{ass:construct}, \ref{ass:budget},
\ref{ass:global}, and \ref{ass:sparsity-jump} hold for Stage~A.  Let
\(r_n\ge r_0\) and let
\[
  N_{r_n}:=\card{\cA_{r_n}}\asymp r_n^d.
\]
Assume
\begin{equation}\label{eq:stage-a-rates}
  \frac{N_{r_n}}{n}\ll\lambda_n\ll\frac1{r_n},
  \qquad
  m_n=o(n/r_n)
  \quad\text{if the minimum-leaf constraint is imposed.}
\end{equation}
Equivalently, existence of such a penalty requires
\[
  r_n^{d+1}=o(n).
\]
Then every Stage-A penalised global minimiser
\[
  \hat t^A
  \in
  \argmin_{\tau\in\cT_{\bar L_A}(\cA_{r_n})}
  \bigl\{R_n(\tau)+\lambda_n\card\tau\bigr\}
\]
has, with probability tending to one, a collapsed split set
\(\topolcoll(\hat t^A)\) in bijection with the true boundaries.  Under this
bijection, each recovered split lies within \(C/r_n\) of its corresponding
true boundary.
\end{proposition}

\begin{proof}[Proof sketch]
The argument pairs one global noise quantity with a boundary-local design
comparison. A feasible comparator $\tau_r^\star$ is built by placing two
consecutive grid splits around every true boundary interior to an atom and one
split at every boundary already on a cutpoint; it represents the atom-projected
truth $\gbar_r$ exactly, and under $m_n=o(n/r_n)$ binomial concentration makes
it Stage-A feasible with probability tending to one. All atom-mean noise is
then absorbed into the single quadratic quantity
$V_r=\sum_a\widehat\omega_aD_a^2$, whose expectation is at most
$N_r\sigma^2/n$, so that $V_r=O_p(N_r/n)=o_p(\lambda_n)=o_p(r^{-1})$ by
Markov's inequality alone --- with no union bound over atoms, no
homoskedasticity, and no global lower bound on atom occupancies. Separating
that noise from occupancy imbalance by means of the empirical-design signal
$\widehat S_r(\tau)$ of~\eqref{eq:empirical-design-signal}, the empirical ANOVA
identity together with contraction of $I-\widehat\Pi_\tau$ yields the sandwich
$R_n(\tau)-R_n(\tau_r^\star)\ge\widehat S_r(\tau)
-2\{\widehat S_r(\tau)V_r\}^{1/2}-V_r$, in which signal, cross term, and noise
are all measured in the same empirical inner product. The population margin of
Lemma~\ref{lem:topol-margin} transfers to $\widehat S_r$ on the finitely many
fixed boundary patches by a local relative Bernstein bound for a
fixed-VC-dimension family of axis-aligned sets, which is the one place the rate
condition $r_n^{d+1}=o(n)$ is used. Two exclusions then finish the argument:
an omitted boundary loses by $\Theta(r_n^{-1})$, more than the largest possible
penalty advantage $\lambda_n\bar L_A=o(r_n^{-1})$, while an unmatched collapsed
split can be contracted without changing the population fit, so its empirical
advantage is at most $4V_r=o_p(\lambda_n)$. The full proof is given in the
Supplementary Material, \S\ref{sec:supp-topology}.
\end{proof}

\paragraph{Proof audit.}
The stochastic argument uses two distinct controls: \(V_r\) controls
response-mean noise globally through Markov's inequality, while the
empirical-design signal is controlled only on boundary-local patches.
The rates combine as
\[
  N_{r_n}/n=o(\lambda_n),
  \qquad
  \lambda_n=o(r_n^{-1}),
  \qquad
  N_{r_n}\asymp r_n^d,
\]
which requires \(r_n^{d+1}=o(n)\).  The feasibility conclusion additionally
uses \(m_n=o(n/r_n)\).  Failure can occur if transition leaves are excluded
by the minimum-leaf constraint, if the local density vanishes, or if the
solver is not globally exact.

\begin{proposition}[Sequential recovery of continuum boundaries]
\label{prop:boundary-recovery}
Under Assumptions~\ref{ass:iid}, \ref{ass:design},
\ref{ass:moments}, \ref{ass:construct}, \ref{ass:budget},
\ref{ass:global}, and \ref{ass:sparsity-jump}, let
\(r_n\ge r_0\) and \(\lambda_n\to0\) satisfy
\[
  \frac{N_{r_n}}{n}\ll\lambda_n\ll\frac1{r_n},
  \qquad
  N_{r_n}\asymp r_n^d,
  \qquad
  m_n=o(n/r_n)
\]
when Stage~A retains the minimum-leaf constraint.  Equivalently, the first
two conditions require \(r_n^{d+1}=o(n)\).  Under the inflated Stage-A
budgets \(\bar L_A=\min\{2\bar L-1,N_{r_n}\}\) and \(2d\),
\[
  \PP\left(
    \begin{array}{c}
    \topolcoll(\hat t^A)\text{ is in bijection with the true boundaries,}\\
    \text{and each recovered boundary is within }C/r_n
    \text{ of its match}
    \end{array}
  \right)
  \longrightarrow1 .
\]
\end{proposition}

\begin{proof}[Proof sketch]
Every hypothesis of Proposition~\ref{prop:parsimony-grid} holds along the joint
sequence $(n,r_n,\lambda_n,m_n)$, so its bijection and $C/r_n$ localisation
apply directly; no fixed-$r$, then $n\to\infty$, interchange is used. The only
additional step converts that localisation into the bracket statement Stage~B
requires: two distinct true boundaries on one coordinate are at distance at
least $\varsigma$ by (Sep), and each recovered representative lies within
$C/r_n$ of its match, so the triangle inequality gives a recovered separation
of at least $\varsigma-2C/r_n\to\varsigma$. Hence the bracket bounded by the
nearest \emph{other} same-coordinate recovered split has width bounded below by
a fixed positive constant and does not shrink at rate $r_n^{-1}$, while still
containing the true boundary. The full proof is given in the Supplementary
Material, \S\ref{sec:supp-topology}.
\end{proof}

\paragraph{Proof audit.}
The conclusion is sequential and follows directly from
Proposition~\ref{prop:parsimony-grid}.  Separation, rather than grid mesh,
sets the eventual Stage-B bracket width.  The bracket conclusion can fail if
same-coordinate true boundaries approach one another with \(n\).

\begin{proposition}[Transition leaves inflate the grid-optimal leaf count]
\label{prop:transition-leaves}
Under Assumption~\ref{ass:sparsity-jump} and for all $r\ge r_0$ (so that no
grid atom contains more than one true boundary), let $L^{\mathrm{cont}}$ be
the leaf count of the true continuum partition and let $B_r$ be the
number of its boundaries not lying at a cutpoint of $\cA_r$. Then
\[
  L^{\mathrm{cont}}\ \le\ L^{\min}_{r}
  \ \le\ L^{\mathrm{cont}}+B_r
  \ \le\ 2L^{\mathrm{cont}}-1 ,
\]
where $L^{\min}_{r}:=\min\{\card\tau:\tau\in\cS^{\cA_r}\}$, and the middle
inequality is an equality whenever the jump values on the two sides of each
such boundary differ from the atom-averaged value between them.
\end{proposition}
\begin{proof}[Proof sketch]
A true boundary interior to an atom $a$ makes $\gbar_r$ take on $a$ an
intermediate value distinct from its values on both neighbouring atoms, by
clause (J) together with the positive mass of both parts of $a$ under (Dens).
By Lemma~\ref{lem:grid-optimal} a grid-optimal tree must therefore separate $a$
from both neighbours, costing exactly one leaf beyond the $L^{\mathrm{cont}}$
leaves of the snapped partition, whereas a boundary already lying at a cutpoint
costs nothing; summing over the at most $L^{\mathrm{cont}}-1$ boundaries gives
the outer bounds, and the middle inequality becomes an equality exactly when
every transition atom's projected value differs from both neighbouring values.
The collapse map $\topolcoll$ then merges each such transition pair back into
the single interval it spans, which is the interval on which Stage~B searches.
The full proof is given in the Supplementary Material,
\S\ref{sec:supp-topology}.
\end{proof}

\paragraph{Proof audit.}
The restriction \(r\ge r_0\) is essential: without it, one atom could contain
multiple true boundaries and the one-extra-leaf-per-boundary count would
fail.  The upper bound is constructive; equality additionally requires each
transition atom to have a projected value distinct from both neighbouring
values.


\section{Superconsistent threshold refinement (Stage B)}\label{sec:stageb}

\begin{remark}[Notation used in this section]
This section introduces several objects distinguished only by
subscript, which are collected here rather than at each first use. Three
distinct $S$-type quantities appear: $S_r(\tau)$, the population margin of
Lemma~\ref{lem:topol-margin}; $\widehat S_r(\tau)$, its
empirical-design analogue, defined
in~\eqref{eq:empirical-design-signal}; and
$S^{\mathrm{act}}_{n,\ell}$/$S^{\mathrm{or}}_{n,\ell}$, the actual- and
oracle-row-set subtree sums of squared error in
Lemma~\ref{lem:ancestor-localisation}'s proof. Two distinct
maximum-order-statistic quantities also appear: $M$, the deterministic
ancestor-collar radius of Lemma~\ref{lem:ancestor-localisation}; and $M_k$,
an intermediate index in the unrelated enumeration of
Proposition~\ref{prop:twostage-superconsistency}'s proof. None of these
symbols denotes the same object as any other, and no result below relies
on conflating them.
\end{remark}

Throughout this section, $\gamma_0$, $Z$, and the design measure $\nu$ are as
fixed in \S\ref{sec:setup-basics}: $Z$ is the response, $\gamma_0(x)=\E[Z\mid
X=x]$ the tree-structured regression function, and $\nu$ the design measure
against which every mass and risk is taken. Multiplying an empirical criterion
by a positive random normalising constant does not change its minimiser; the
refinement problem is therefore a squared-error change-point problem under the
design measure \(\nu\).  Assumption~\ref{ass:moments} supplies
\[
  \sup_x \E\!\left[
    \{Z-\gamma_0(X)\}^2\mid X=x
  \right]<\infty ,
\]
which is the only distributional property of $Z$ used below beyond
Assumption~\ref{ass:iid}.

Throughout this section, set
\[
  p_{\min}:=f_{\min}\,c_{\mathrm{box}}\,q_{\min},
\]
with \(f_{\min}\) the density lower bound of clause (Dens), \(q_{\min}\) the
node-mass floor of clause (Node), and \(c_{\mathrm{box}}\) the
cross-sectional constant fixed before Lemma~\ref{lem:topol-margin}. The
three factors are not interchangeable, and the \((d-1)\)-dimensional factor
cannot be omitted: \(\omega_{N,r}(c)\) is the \(\nu\)-mass of a
\(d\)-dimensional prism cell of thickness \(\abs{c-c_0}\), so it carries the
cell's \((d-1)\)-dimensional cross-sectional size as well as a density.
Explicitly, writing \(G_r\) for the \((d-1)\)-dimensional footprint of prism
cell \(r\) in the coordinates other than \(q\), clause (Dens) gives
\[
  \omega_{N,r}(c)
  \ \ge\
  f_{\min}\,\operatorname{Leb}_{d-1}(G_r)\,\abs{c-c_0}
  +o\!\left(\abs{c-c_0}\right).
\]
The footprints \(G_r\) are disjoint and, by Lemma~\ref{lem:prism-profile}(i)--(ii),
exhaust the cross-section \(\Gamma_N\) of the box \(R_N\) transverse to
\(q\); under (Local-Jump) every resulting cell is active. Summing therefore
gives exactly
\[
  \sum_r\omega_{N,r}(c)
  \ \ge\
  f_{\min}\operatorname{Leb}_{d-1}(\Gamma_N)\abs{c-c_0}
  +o\!\left(\abs{c-c_0}\right),
\]
and, as fixed above, \(\operatorname{Leb}_{d-1}(\Gamma_N)\ge c_{\mathrm{box}}
q_{\min}\). Absorbing the \(o(\abs{c-c_0})\) term by shrinking the
neighbourhood yields the bound
\[
  \sum_r\omega_{N,r}(c)\ \ge\ \frac{p_{\min}}{2}\abs{c-c_0}
\]
used below. Making the cross-sectional factor explicit weakens every drift
constant below by the fixed positive factor \(c_{\mathrm{box}}q_{\min}\),
and changes nothing else: no condition in this paper requires
\(p_{\min}\) to exceed a specified threshold.

\subsection{Consistency and reduction to a fixed bracket}
\label{sec:stageb-bracket}

Both tiers below face the same preliminary difficulty before any rate
argument applies: the bracket $\widehat B_n$ that Stage B actually searches
(Assumption~\ref{ass:solver-twostage}(e)) is a random set, produced by
Stage A, whereas the classical results this section relies on
(\citealt{chan1993,seijosen2011}, and Lemma~\ref{lem:prism-profile}'s own
localisation argument below) presume a fixed, data-independent search
region. The following lemma removes this randomness once and for all,
before either tier is entered.

\begin{lemma}[Consistency and reduction to a fixed bracket]
\label{lem:bracket-reduction}
Fix an internal node $N$ on coordinate $q$ with true threshold
$c_0:=c_{0,N}$, and let $[q_N^-,q_N^+]$ denote the fixed, compact extent of
the true node box $R_N$ along coordinate $q$ (well defined and bounded, by
the construction of $R_N^+$ and the bounded coordinate diameters assumed in
clause (Dens)). (Here and below, $R_N$ and $R_N^+$ always denote the node's
axis-aligned \emph{box}, per clause (Dens); a superscripted $R_N^{\bullet}$
with $\bullet\in\{X,\mathrm{stump},\mathrm{prof}\}$ always denotes a
real-valued \emph{risk}. The two uses do not otherwise overlap.) Write
$R_N^X$ for the relevant population risk as a
function of the candidate threshold --- $R_N^{\mathrm{stump}}$ if both
children of $N$ are terminal, $R_N^{\mathrm{prof}}$ otherwise --- and let
$\widehat c_N$ denote the Stage-B minimiser of the corresponding empirical
criterion over $\widehat B_n$. Let
\[
  B^-:=\bigl(c_0-\varsigma/4,\,c_0+\varsigma/4\bigr),
\]
and let $\widetilde c_N$ denote the minimiser of the same empirical
criterion over the FIXED interval $B^-$. Then:
\begin{enumerate}[label=(\alph*)]
  \item (Global separation) For every $\eta\in(0,q_N^+-q_N^-]$,
  \[
    \delta_N(\eta)
    :=
    \inf\Bigl\{R_N^X(c)-R_N^X(c_0):\
      c\in[q_N^-,q_N^+],\ \abs{c-c_0}\ge\eta\Bigr\}
    >0.
  \]
  In particular $\delta_N:=\delta_N(\varsigma/4)>0$ is the value used in
  part (c).
  \item (Consistency) $\widehat c_N\to_p c_0$.
  \item (Bracket reduction) $\PP(\widehat c_N=\widetilde c_N)\to1$.
\end{enumerate}
\end{lemma}

\begin{proof}[Proof sketch]
Part (a) is population-only: the ANOVA identity makes
$R_N^X(c)-R_N^X(c_0)$ nonnegative, and strictly positive for $c\ne c_0$
because the slab $\mathcal I_N(c,c_0)$ has strictly positive $\nu$-mass by
(Dens) while (Local-Jump) gives one of its cells a nonzero jump, so $\tau(c)$
is not sufficient; continuity of $R_N^X$ plus compactness of
$\{c\in[q_N^-,q_N^+]:\abs{c-c_0}\ge\eta\}$ then forces a strictly positive
infimum. Part (b) applies the standard argmin-consistency argument for
M-estimators on a fixed compact parameter space \citep{vandervaart1998} to an
auxiliary minimiser over the deterministic set $[q_N^-,q_N^+]$, the required
uniform convergence following from Glivenko--Cantelli applied to the finitely
many partial sums indexed by the monotone indicator $\mathbf 1\{X_q\le c\}$,
together with Assumption~\ref{ass:moments}'s second-moment bound. Part (c) is
where the random bracket is discharged: Proposition~\ref{prop:boundary-recovery}
gives $\PP(B^-\subseteq\widehat B_n)\to1$ while
$\widehat B_n\subseteq[q_N^-,q_N^+]$ holds always, and a global minimiser over
the largest of three nested sets that lands in the smallest is automatically
the minimiser over each of them, which identifies the auxiliary minimiser with
$\widehat c_N$ and with $\widetilde c_N$ simultaneously. The full proof is
given in the Supplementary Material, \S\ref{sec:supp-stageb}.
\end{proof}

\paragraph{Proof audit.} Part (a) uses only (Dens), (Local-Jump), and
compactness --- no probabilistic tool. Part (b) uses only a
coordinatewise Glivenko--Cantelli argument on a FIXED compact domain, not
the random $\widehat B_n$. Part (c) is where the random bracket is
discharged, using Proposition~\ref{prop:boundary-recovery}'s
already-established width guarantee; nothing here duplicates that
proposition's own proof. The lemma can fail if $\varsigma$ (the separation
constant) is not fixed, or if the node box $R_N$ is unbounded.

Lemma~\ref{lem:bracket-reduction} and Lemma~\ref{lem:prism-profile} below are
both stated with every descendant threshold at its \emph{true} value.
Stage~B does not have access to true descendant values: by
Assumption~\ref{ass:solver-twostage}(b), a node's refinement proceeds
level-by-level with every strictly-deeper threshold still held at its
Stage-A grid value. The next lemma closes this gap.

\begin{lemma}[Stability under Stage-A descendant thresholds]
\label{lem:descendant-stability}
Let \(N\) be an internal node and let \(\widetilde R^X_{N,n}\) denote the
population criterion of Lemma~\ref{lem:bracket-reduction} computed with
every descendant threshold of \(N\) held at its Stage-A value rather than at
its true value. On the event \(E_n\) of
Proposition~\ref{prop:boundary-recovery}, every such descendant threshold
lies within \(C/r_n\) of its true value. Then, with probability tending to
one:
\begin{enumerate}[label=(\alph*),nosep]
  \item \(\sup_{c\in[q_N^-,q_N^+]}
    \abs{\widetilde R^X_{N,n}(c)-R_N^X(c)}=O(1/r_n)=o(1)\);
  \item for every fixed \(\eta>0\),
    \(\inf_{\abs{c-c_0}\ge\eta}
      \{\widetilde R^X_{N,n}(c)-\widetilde R^X_{N,n}(c_0)\}
      \ge\tfrac12\delta_N(\eta)>0\);
  \item the perturbed prism masses and jump marks of
    Lemma~\ref{lem:prism-profile} satisfy
    \(\widetilde\omega_{N,r}(c)=\omega_{N,r}(c)\{1+O(1/r_n)\}\) and
    \(\abs{\widetilde\Delta_{N,r}}\ge\kappa_{\min}/2\), so
    \eqref{eq:prism-drift} holds for \(\widetilde R^X_{N,n}\) with
    \(p_{\min}\) replaced by \(p_{\min}/2\).
\end{enumerate}
Consequently Lemmas~\ref{lem:bracket-reduction} and~\ref{lem:prism-profile}
apply verbatim to the criterion Stage~B actually minimises, with the drift
constant halved.
\end{lemma}
\begin{proof}[Proof sketch]
On the event $E_n$ every descendant threshold is displaced by $O(1/r_n)$,
which by the upper density bound of (Dens) and the bounded coordinate
diameters perturbs the $\nu$-mass of the affected leaf regions by $O(1/r_n)$;
part (a) follows from this and the boundedness of the profiled leaf means
(Assumptions~\ref{ass:design} and~\ref{ass:construct}(c)). Part (b) combines
(a) with Lemma~\ref{lem:bracket-reduction}(a): the perturbed infimum lies
within $O(1/r_n)$ of the true one and so exceeds $\delta_N(\eta)/2$ once
$1/r_n<\delta_N(\eta)/2$. Part (c) follows because each perturbed
adjacent-leaf pair differs from a true pair only through regions of $\nu$-mass
$O(1/r_n)$, whose profiled means therefore differ by $O(1/r_n)$ given the
leaf-mass floor of (Node), so (Local-Jump) transfers to the perturbed pairs
with $\kappa_{\min}$ replaced by $\kappa_{\min}/2$. The full proof is given in
the Supplementary Material, \S\ref{sec:supp-stageb}.
\end{proof}

\subsection{Two tiers of refinement}\label{sec:stageb-tiers}

The proof has two tiers. Both first invoke
Lemma~\ref{lem:bracket-reduction} to replace the random bracket $\widehat
B_n$ by the fixed interval $B^-$; the tiers differ only in which criterion
--- stump or profiled --- is then analysed on that fixed interval.

\paragraph{Tier 1: leaf-parent nodes.}
If both children of \(N\) are terminal, the Stage-B objective is exactly the
classical two-group stump criterion, taken under the design measure \(\nu\)
rather than under \(P_X\); since \(\nu\) is a fixed measure equivalent to
\(P_X\) by Assumption~\ref{ass:design}, and the observations are i.i.d.\ by
Assumption~\ref{ass:iid}, the change-point results below apply unchanged with
\(\nu\) in place of \(P_X\) throughout.
By Lemma~\ref{lem:bracket-reduction}(c),
the actual Stage-B minimiser $\widehat c_N$ equals $\widetilde c_N$, the
minimiser of the same criterion over the FIXED interval $B^-$, with
probability tending to one. The random-design change-point results of
\citet{chan1993} and \citet{seijosen2011}, whose
conditions are supplied here by (Dens), (Node), (Local-Jump), and the standing
second-moment condition, apply directly to $\widetilde c_N$ (a legitimate,
unconditionally-defined statistic on the fixed set $B^-$) and give
\[
  \limsup_{n\to\infty}
  \PP\!\left(n\abs{\widetilde c_N-c_{0,N}}>u\right)
  =:\pi_N(u),
  \qquad
  \pi_N(u)\longrightarrow0
  \quad\text{as }u\to\infty .
\]
Since $\PP(\widehat c_N=\widetilde c_N)\to1$, this transfers to $\widehat
c_N$:
\[
  n(\widehat c_N-c_{0,N})=O_p(1).
\]

\paragraph{Tier 2: non-terminal children.}
If at least one child is non-terminal, replacing each child subtree by a
single group is generally invalid.  In particular, a descendant split on the
same coordinate can make the resulting two-group population criterion attain
its minimum away from \(c_{0,N}\).  The Stage-B criterion must instead retain
and profile all descendant leaves.  The following lemma treats precisely that
criterion.

\begin{lemma}[Prism decomposition of the local profiled criterion]
\label{lem:prism-profile}
Fix an internal node \(N\), its splitting coordinate \(q\), its true threshold
\(c_0:=c_{0,N}\), and its true region \(R_N\).  Suppose that the topology,
the true ancestor row set, and the population-level descendant cell structure
are fixed.  (The corresponding statement with descendant thresholds held at
their Stage-A values instead of their true values follows from
Lemma~\ref{lem:descendant-stability}, with \(p_{\min}\) replaced by
\(p_{\min}/2\) throughout.)  For a candidate \(c\) in the node's true box (in particular for
\(c\) in the fixed interval \(B^-\) of Lemma~\ref{lem:bracket-reduction},
which is where parts (iv)--(v) below are ultimately applied), let
\[
  \mathcal I_N(c,c_0)
  :=
  R_N\cap
  \bigl\{x:x_q\in(c\wedge c_0,c\vee c_0]\bigr\}.
\]
Then:

\begin{enumerate}[label=(\roman*),leftmargin=*]
  \item \(\mathcal I_N(c,c_0)\) is exactly the set whose side-of-\(N\)
  membership differs under \(c\) and \(c_0\).

  \item The slab \(\mathcal I_N(c,c_0)\) admits a disjoint decomposition into
  at most
  \[
    L_N^-L_N^+
    \le \frac{(L_N^-+L_N^+)^2}{4}
    \le \frac{\bar L^2}{4}
  \]
  active prism cells, where \(L_N^-\) and \(L_N^+\) are the numbers of
  descendant leaves below the two children of \(N\).

  \item If \(\omega_{N,r}(c)\) is the \(\nu\)-mass of prism cell \(r\), and
  \(\Delta_{N,r}\) is the jump between its adjacent leaves, then for every
  cell \(r\) whose \((d-1)\)-dimensional footprint \(G_r\) has
  \(\operatorname{Leb}_{d-1}(G_r)>0\),
  \[
    \omega_{N,r}(c)=\Theta\!\left(\abs{c-c_0}\right),
    \qquad
    \kappa_{\min}\le\abs{\Delta_{N,r}}\le\kappa_{\max}.
  \]
  A cell with \(\operatorname{Leb}_{d-1}(G_r)=0\) (a source--target leaf
  pair whose footprints do not overlap) has \(\omega_{N,r}\equiv0\) and
  contributes zero to both the cell count of (ii) and the drift sum of (iv);
  it is retained in the count only as an upper-bound artefact.
  Thus, at \(c=c_0+u/n\) with fixed \(u\), every cell with positive
  footprint has \(\omega_{N,r}(c)=\Theta(n^{-1})\).

  \item Writing \(R_N^{\mathrm{prof}}(c)\) for the population squared-error
  risk after profiling all descendant-leaf means,
  \begin{equation}
  \label{eq:prism-risk}
    R_N^{\mathrm{prof}}(c)-R_N^{\mathrm{prof}}(c_0)
    =
    \sum_r \omega_{N,r}(c)\Delta_{N,r}^2
    +O\!\left(\abs{c-c_0}^2\right).
  \end{equation}
  Consequently, for all sufficiently small \(\abs{c-c_0}\),
  \begin{equation}
  \label{eq:prism-drift}
    R_N^{\mathrm{prof}}(c)-R_N^{\mathrm{prof}}(c_0)
    \ge
    \frac{\kappa_{\min}^2p_{\min}}{4}
    \abs{c-c_0}.
  \end{equation}

  \item Let \(\widetilde c_N\) minimise the corresponding empirical,
  fully-profiled subtree SSE over observed values in the FIXED interval
  \(B^-\) of Lemma~\ref{lem:bracket-reduction}, and let \(\widehat c_N\) be
  the actual Stage-B output (the same minimisation over the recovered
  bracket \(\widehat B_n\)). Then
  \[
    n(\widetilde c_N-c_0)=O_p(1),
    \qquad\text{and hence, by Lemma~\ref{lem:bracket-reduction}(c),}\qquad
    n(\widehat c_N-c_0)=O_p(1).
  \]
\end{enumerate}
\end{lemma}

\begin{proof}[Proof sketch]
The argument has four steps. Step~1 identifies the moving slab: only an
observation whose $q$-coordinate lies between $c$ and $c_0$ changes side of
$N$, which is (i). Step~2 refines that slab into at most
$L_N^-L_N^+\le\bar L^2/4$ disjoint prism cells indexed by realised
source--target leaf pairs; this is legitimate because
Assumption~\ref{ass:solver-twostage}(b)'s bracket together with (Sep) prevents
a candidate from crossing any other $q$-split, so every same-coordinate
descendant test is constant throughout the slab and may be absorbed into a
fixed path label, while cell masses and jump marks are read off (Dens) and
(Local-Jump) --- giving (ii) and (iii). Step~3 applies the weighted ANOVA
identity one destination leaf at a time; the pooling correction is
$O(\abs{c-c_0}^2)$ once (Node) bounds the correctly assigned core mass below,
which yields~\eqref{eq:prism-risk} and then, with the cross-sectional mass
bound $\sum_r\omega_{N,r}(c)\ge p_{\min}\abs{c-c_0}/2$, the linear drift
bound~\eqref{eq:prism-drift}. Step~4 rescales to $c=c_0+u/n$, where every
active cell holds $\Theta(1)$ observations, and represents the local
unnormalised criterion as a finite sum --- at most $\bar L^2/4$ terms, a number
fixed in $n$ --- of marked rare-event counting processes with aggregate
one-sided drift at least $\kappa_{\min}^2p_{\min}\abs u/2$; this is the point
at which a tree differs from a two-group problem, since no single change-point
process carries the drift. Consistency from
Lemma~\ref{lem:bracket-reduction}(b) confines the minimiser to shrinking
neighbourhoods of $c_0$, on which the argmax theorem for discontinuous,
locally-defined criteria \citep{seootsu2018} applies and yields tightness of
$n(\widetilde c_N-c_0)$; Lemma~\ref{lem:bracket-reduction}(c) transfers it to
the actual Stage-B output, proving (v). The full proof is given in the
Supplementary Material, \S\ref{sec:supp-stageb}.
\end{proof}

\paragraph{Proof audit.}
All structural conditions used above are contained in
Assumption~\ref{ass:sparsity-jump}; the fixed bracket and exact minimisation
come from Assumption~\ref{ass:solver-twostage}.  The cited change-point
results are used for their rate-\(n\) localisation machinery, while the new
prism and joint-process arguments verify the feature not covered directly by
those references.  The linear drift constant is explicitly
\(\kappa_{\min}^2p_{\min}/4\), and the quadratic ANOVA remainder cannot change
the rate.  The argument can fail if an active adjacent pair has zero jump, if
the local design density vanishes, if leaf masses approach zero, or if the
number of leaves grows with \(n\).  With fixed nonzero jumps, the \(n^{-1}\)
rate is the usual sharp change-point rate.

\begin{remark}[Profiling the descendant-leaf means]
\label{rem:profiling-correction}
Lemma~\ref{lem:prism-profile} remains valid when every descendant-leaf mean is
re-fitted at every candidate threshold, as Stage B requires.  To see this
explicitly, fix \(\abs u\le U\), put \(h=u/n\), and let \(L(c)\) be a candidate
leaf.  Since its population mass is at least \(q_{\min}/2\), its candidate
population mean satisfies
\[
  \mu_{L(c)}-\mu_{L(c_0)}
  =
  O(\abs h).
\]
Indeed, only \(O(n\abs h)\) observations in expectation are reassigned, their
population contribution is \(O(\abs h)\), and the leaf denominator is bounded
away from zero.  Conditional sampling fluctuation contributes the usual
root-\(n\) term, so uniformly over the finitely many leaves and
\(\abs u\le U\),
\begin{align*}
  \widehat\mu_{L(c)}-\mu_{L(c_0)}
  &=
  \{\mu_{L(c)}-\mu_{L(c_0)}\}
  +\{\widehat\mu_{L(c)}-\mu_{L(c)}\}
  \\
  &=
  O(\abs{c-c_0})+O_p(n^{-1/2}).
\end{align*}
The first term is deterministic misclassification bias; the second is
sampling noise under a fixed leaf assignment.

Because squared error is quadratic in the leaf mean, the resulting
population-risk correction is
\[
  O\!\left(\abs{c-c_0}^2\right)
  +
  O_p\!\left(\abs{c-c_0}\,n^{-1/2}\right).
\]
More precisely, if \(N_L(c)\) is the leaf count and
\(d_L(c)=\widehat\mu_{L(c)}-\mu_{L(c)}\), the least-squares identity gives
\[
  \sum_{i\in L(c)}\{Z_i-\widehat\mu_{L(c)}\}^2
  =
  \sum_{i\in L(c)}\{Z_i-\mu_{L(c)}\}^2
  -
  N_L(c)d_L(c)^2.
\]
For \(\abs u\le U\), only \(O_p(1)\) observations change assignment, whence
\[
  N_L(c)-N_L(c_0)=O_p(1),
  \qquad
  d_L(c_0)=O_p(n^{-1/2}),
  \qquad
  d_L(c)-d_L(c_0)=o_p(n^{-1/2}).
\]
The last display follows because only \(O_p(1)\) of the leaf's \(\Theta(n)\)
observations are reassigned, contributing a mean shift of
\(O_p(1)/\Theta(n)=O_p(n^{-1})\), strictly smaller than the leaf's own
\(O_p(n^{-1/2})\) sampling fluctuation.
Therefore,
\begin{align*}
  &N_L(c)d_L(c)^2-N_L(c_0)d_L(c_0)^2
  \\
  &\quad=
  \{N_L(c)-N_L(c_0)\}d_L(c_0)^2
  +N_L(c)\{d_L(c)-d_L(c_0)\}
       \{d_L(c)+d_L(c_0)\}
  \\
  &\quad=
  O_p(n^{-1})+o_p(n^{-1/2})O_p(n^{1/2})
  =
  o_p(1)
\end{align*}
on the unnormalised local-SSE scale, uniformly over \(\abs u\le U\).
On the normalised risk scale, the deterministic and cross terms are
\(O(n^{-2})\) and \(O_p(n^{-3/2})\), whereas
\eqref{eq:prism-drift} is at least
\(\kappa_{\min}^2p_{\min}\abs u/(4n)\).  Thus profiling contributes only a
second-order perturbation and does not change the \(n^{-1}\) threshold rate.
The conclusion concerns threshold locations; the profiled leaf means
themselves retain their ordinary \(O_p(n^{-1/2})\) rate.
\end{remark}

\subsection{Ancestor-refinement contamination}\label{sec:stageb-ancestor}

\begin{lemma}[Contamination-robust levelwise localisation]
\label{lem:ancestor-localisation}
Suppose the conclusions of Lemma~\ref{lem:prism-profile} and
Remark~\ref{rem:profiling-correction} hold for a node at level \(\ell\) when
its row set is defined by the true ancestor thresholds.  If every ancestor
threshold has already been shown to satisfy
\[
  n(\widehat c_a-c_{0,a})=O_p(1),
\]
then the same conclusion holds when the level-\(\ell\) Stage-B criterion uses
the row set induced by the refined ancestor thresholds:
\[
  n(\widehat c_\ell-c_{0,\ell})=O_p(1).
\]
\end{lemma}

\begin{proof}[Proof sketch]
The induction hypothesis is used only to place every ancestor error inside a
deterministic collar $\mathcal C_{n,M}$ of half-width $M/n$, at the cost of one
third of an $\epsilon$-budget. Clause (Joint) and a bare union bound then show
that, with probability tending to one, no observation lies simultaneously in
that collar and in the level-$\ell$ window
$\abs{x_{q_\ell}-c_{0,\ell}}\le U/n$, the relevant probability being
$O(MU/n)$; when ancestor and descendant share a coordinate, (Sep) makes the two
collars disjoint outright. On that event a contaminated observation's
prism-cell membership does not vary as $c$ ranges over the local window, so,
using the $\Theta_p(n)$ leaf denominators granted by
Assumption~\ref{ass:solver-twostage}(g), it perturbs the actual-minus-oracle
local contrast by only $O_p(n^{-1})=o_p(1)$ on the unnormalised SSE scale.
Because that control holds on a fixed compact $u$-window, escape must be
excluded separately: every contributing observation lies in
$\mathcal C_{n,M}$ \emph{regardless} of $u$ --- collar membership is a property
of the ancestor coordinate alone --- so the discrepancy is $O_p(1)$ uniformly
over $\abs u\le\varsigma n/4$ with no peeling and no $u$-indexed probability
calculation, and the escape-exclusion step that underlies the oracle
criterion's own tightness (\citealt{chan1993,seijosen2011} in Tier~1,
\citealt{seootsu2018} for the profiled criterion) then forces the actual
criterion to exceed its value at $c_{0,\ell}$ throughout
$U\le\abs u\le\varsigma n/4$. Combining the two ranges gives
$n(\widehat c_\ell-c_{0,\ell})=O_p(1)$. The full proof is given in the
Supplementary Material, \S\ref{sec:supp-stageb}.
\end{proof}

\paragraph{Proof audit.}
The induction hypothesis is used only to construct the deterministic
ancestor collars.  Clause (Joint) is used only when ancestor and descendant
coordinates differ; (Sep) handles repeated coordinates.  The collar contains
\(O_p(1)\), not deterministically bounded, observations.  Direct
differentiation gives an \(O_p(n^{-1})=o_p(1)\) perturbation of the
unnormalised local contrast, which is dominated by its \(O_p(1)\) marked
change-point fluctuations and \(O(1)\) drift.  The argument would fail if the
joint density had a singular component concentrated at the crossing of two
boundaries, or if node masses vanished with \(n\).

\subsection{All levels and all recovered boundaries}\label{sec:stageb-all}

At the root there is no ancestor contamination.  A root whose children are
terminal is covered by Tier~1; otherwise it is covered by
Lemma~\ref{lem:prism-profile}.  Suppose inductively that every boundary above
level \(\ell\) has the \(n^{-1}\) rate.  Lemma~\ref{lem:ancestor-localisation}
then transfers the oracle-row-set rate to every node at level \(\ell\).
Proposition~\ref{prop:boundary-recovery} also places every other recovered
split within $C/r_n$ of its true split.  Since \(r_n\to\infty\), these
other split locations converge to their true values.  Their induced prism
masses, jump marks, and positive drift constants therefore differ from their
population counterparts by \(o_p(1)\), by Lemma~\ref{lem:descendant-stability};
the lower bound in
\eqref{eq:prism-drift} remains bounded away from zero (at half its stated
value) with probability tending
to one.  No rate faster than consistency is required for these non-local
descendant splits.

The induction stops after at most \(d\) levels.  Both \(d\) and the total
number of boundaries, at most \(\bar L-1\), are fixed.

\begin{proposition}[Simultaneous superconsistency of all refined boundaries]
\label{prop:twostage-superconsistency}
Suppose Assumptions~\ref{ass:iid}, \ref{ass:design}, and~\ref{ass:moments}
hold.
Under Assumptions~\ref{ass:sparsity-jump}, including (Local-Jump) and (Joint),
and~\ref{ass:solver-twostage}, and on the topology-recovery conclusion of
Proposition~\ref{prop:boundary-recovery}, every Stage-B-refined boundary
satisfies
\[
  \widehat c_k-c_{0,k}=O_p(n^{-1}),
  \qquad k=1,\ldots,K,
  \qquad K\le\bar L-1.
\]
Equivalently, the result holds simultaneously:
\[
  \max_{1\le k\le K}
  n\abs{\widehat c_k-c_{0,k}}
  =
  O_p(1).
\]
Thus the conclusion applies to an arbitrary fixed-\(\bar L\), fixed-depth
tree, not merely to a two- or three-leaf special case.
\end{proposition}

\begin{proof}[Proof sketch]
The root carries no ancestor contamination, so its rate is the classical stump
argument when both children are terminal and Lemma~\ref{lem:prism-profile}
otherwise; Lemma~\ref{lem:ancestor-localisation} then propagates the rate one
level at a time, and the induction terminates after at most $d$ levels because
the depth is fixed. Simultaneity is a finite union bound over the at most
$\bar L-1$ boundaries: choosing $M_k$ with
$\limsup_n\PP(n\abs{\widehat c_k-c_{0,k}}>M_k)<\epsilon/K$ for each $k$ and
taking $M=\max_{k\le K}M_k$ bounds the maximum. Since $K\le\bar L-1<\infty$, no
growing-dimensional union bound is required. The full proof is given in the
Supplementary Material, \S\ref{sec:supp-stageb}.
\end{proof}

\paragraph{Proof audit.}
The proof uses exact topology recovery, fixed depth, fixed leaf count, exact
Stage-B optimisation, positive leaf masses, nonvanishing adjacent jumps, and
the local density conditions stated before the proposition.  Tier~1 is the
classical single-change-point case; Tier~2 is supplied by the new prism and
profiling arguments rather than by treating a non-terminal subtree as a
stump.  Every local process has \(O(1)\) drift and fluctuation on the
\(n^{-1}\) threshold scale, and all profiling and ancestor-row-set
perturbations are \(o_p(1)\) on the unnormalised SSE scale.  The rate is sharp
for discontinuous regression boundaries.  Failure is localised to vanishing
jumps, vanishing face density, collapsing leaf mass, singular
ancestor--descendant designs, imperfect optimisation, or a leaf count or
depth that grows with \(n\).


\section{Implementation}\label{sec:implementation}
[TODO -- not yet drafted; scope revised 2026-09-01 for JASA Theory \&
Methods conventions (see
\texttt{quality\_reports/plans/2026-09-01\_two-paper-split.md} \S7): JASA
treats software as a reproducibility/supplementary-materials artifact, not
a primary main-text section. This should be a SHORT paragraph (not a
systems-style section) -- name the \texttt{optimaltrees} package and its
\texttt{RefinedTreeModel} S7 class, \texttt{as\_coordinate\_tree()},
\texttt{collapse\_transitions()}, \texttt{refine\_tree\_cuts()} by
function, point to the package's own documentation for engineering
detail, and close with a forward pointer to the Supplementary Material
section (\S\ref{sec:supplement}) for the reproducibility statement. Do
not expand this into a full package description -- that belongs in
\texttt{optimaltrees}'s own documentation, which already exists.]

\section{Simulations}\label{sec:simulations}
[TODO -- not started. Deferred per the 2026-09-01 planning discussion until
after both papers' outlines are drafted. Candidate content: threshold-MSE
scaling ($\propto n^{-2}$), topology-recovery frequency vs.\ grid
resolution $r$, comparison against greedy CART and grid-only optimal trees,
timing.]

\section{Discussion}\label{sec:discussion}
[TODO.]

% ===========================================================================
% Supplementary Material. Unnumbered (starred) headings, per JASA's convention
% that the supplement is a block distinct from the numbered main text. The
% section/subsection counters are stepped by hand so that \ref{} to these
% headings still resolves to the S / S.1 / S.2 / S.3 tags printed with them,
% and \theHsection/\theHsubsection are redefined so hyperref's anchors do not
% collide with those of the main text's Section~1 and its subsections.
\renewcommand{\thesection}{S}
\renewcommand{\theHsection}{supplement.\arabic{section}}
\renewcommand{\thesubsection}{S.\arabic{subsection}}
\renewcommand{\theHsubsection}{supplement.\arabic{subsection}}
\setcounter{section}{0}
\setcounter{subsection}{0}

\refstepcounter{section}
\section*{Supplementary Material}\label{sec:supplement}
\addcontentsline{toc}{section}{Supplementary Material}

This supplement collects the full proofs of the results stated in
Sections~\ref{sec:setup}--\ref{sec:stageb}, organised by the section in which
each result appears; the main text retains every statement in place, together
with a proof sketch and a pointer to the subsection below in which its proof is
given. No new assumption, definition, or result is introduced here, and
nothing below is used in the main text beyond the results already stated
there. The two-stage algorithm whose properties are proved below is
implemented and tested in the \texttt{optimaltrees} \textsf{R} package: its
\texttt{RefinedTreeModel} class carries the refined tree, while
\texttt{as\_coordinate\_tree()}, \texttt{collapse\_transitions()}, and
\texttt{refine\_tree\_cuts()} implement respectively the extraction of a
Stage-A topology, the collapse map $\topolcoll$ of
\S\ref{sec:setup-twostage}, and Stage~B's bracketed off-grid scan. Code
sufficient to reproduce every numerical result reported in this paper is
distributed with that package [TODO: add repository URL].


\refstepcounter{subsection}
\subsection*{S.1\quad Proofs for Section~\ref{sec:setup}}\label{sec:supp-setup}
\addcontentsline{toc}{subsection}{S.1\quad Proofs for
  Section~\protect\ref{sec:setup}}

Section~\ref{sec:setup} states the standing assumptions, the sufficient-class
definition, and the two-stage construction; none of its numbered results
carries a proof obligation, so no proof is collected here. The subsection is
retained so that the supplement's organisation matches the main text's
section by section.

\refstepcounter{subsection}
\subsection*{S.2\quad Proofs for Section~\ref{sec:topology}}%
\label{sec:supp-topology}
\addcontentsline{toc}{subsection}{S.2\quad Proofs for
  Section~\protect\ref{sec:topology}}

The proof of Lemma~\ref{lem:grid-optimal} occupies only a few lines and is
retained in place in the main text; the remaining proofs of
\S\ref{sec:topology} follow, in the order in which their results are stated.

\begin{proof}[Proof of Lemma~\ref{lem:topol-margin}]
\emph{Proof strategy.}
We first lower-bound the approximation loss caused by omitting one true
boundary, using pairs of points that differ only in the splitting coordinate.
We then construct a feasible grid-aligned comparator by snapping every true
boundary to a neighbouring cutpoint.  Subtracting the comparator's
\(O(r^{-1})\) approximation loss gives the stated margin.

Fix an omitted true boundary \(b=(q,c_{0,N})\) in the true node \(R_N\), and
write
\[
  h_r:=\frac{C}{2r}.
\]
For each point \(z\) on an active \((d-1)\)-dimensional face cell and each
\(t\in(0,h_r]\), define
\[
  x^-(z,t):=(z,c_{0,N}-t),
  \qquad
  x^+(z,t):=(z,c_{0,N}+t),
\]
with the splitting coordinate inserted in the \(q\)-th position.
By the choice of \(r_0\), these points lie in the true node's local
neighbourhood.  Clause (Dens) gives
\[
  f_{\nu,N}\{x^-(z,t)\}\ge f_{\min},
  \qquad
  f_{\nu,N}\{x^+(z,t)\}\ge f_{\min}
\]
in density notation.  In particular, both sides belong to the support; the
disjoint-support counterexample is therefore excluded.

The two points have identical coordinates other than \(q\).  Hence no split
on another coordinate can separate them.  Because \(\tau\) omits \(b\), no
\(q\)-split of \(\tau\) lies between \(c_{0,N}-h_r\) and
\(c_{0,N}+h_r\): any such raw split would survive as a singleton under
\(\topolcoll\), or would belong to a collapsed transition interval meeting
that range.  Thus \(x^-(z,t)\) and \(x^+(z,t)\) lie in the same leaf of
\(\tau\).

Let \(g_\tau\) be any function constant on that leaf.  By (Local-Jump),
\[
  \abs{\gamma_0\{x^+(z,t)\}
       -\gamma_0\{x^-(z,t)\}}
  \ge\kappa_{\min}.
\]
The triangle inequality implies
\[
  \kappa_{\min}
  \le
  \abs{\gamma_0(x^+)-g_\tau(x^+)}
  +
  \abs{g_\tau(x^-)-\gamma_0(x^-)}.
\]
Therefore at least one summand is at least \(\kappa_{\min}/2\), and hence
\[
  \{\gamma_0(x^+)-g_\tau(x^+)\}^2
  +
  \{\gamma_0(x^-)-g_\tau(x^-)\}^2
  \ge \frac{\kappa_{\min}^2}{4}.
\]
Integrating over \(t\in(0,h_r]\), over all active face cells, and using
(Dens), (Node), and the cross-sectional lower bound gives
\begin{align*}
  R(\tau)-R(\gamma_0)
  &\ge
  \frac{\kappa_{\min}^2}{4}
  f_{\min}
  \bigl(c_{\mathrm{box}}q_{\min}\bigr)h_r\\
  &=
  \frac{c_{\mathrm{box}}}{8}
  \kappa_{\min}^2f_{\min}q_{\min}\frac{C}{r}\\
  &=
  c_{\mathrm{topol}}\kappa_{\min}^2f_{\min}q_{\min}\frac{C}{r}.
\end{align*}
Grid cells that meet the boundary of the fixed face patch can be discarded.
Their partial-overlap volume is \(o(r^{-1})\), because the face patch has
fixed positive \((d-1)\)-dimensional measure and the grid mesh tends to zero.
After increasing \(r_0\), the displayed bound holds with \(c_{\mathrm{box}}\)
replaced by \(c_{\mathrm{box}}/2\), i.e.\ with \(c_{\mathrm{topol}}\) above
replaced by \(c_{\mathrm{topol}}/2\) in the final display, since
\(c_{\mathrm{topol}}:=c_{\mathrm{box}}/8\) is linear in \(c_{\mathrm{box}}\);
\eqref{eq:C-margin-choice}'s choice of \(C\) is made to accommodate this
factor of \(2\) loss, so no circularity arises from fixing \(c_{\mathrm{box}}\) net of
this discard loss before choosing \(C\).

It remains to compare this loss with the best grid tree.  Snap every true
boundary to a nearest grid cutpoint while preserving the true ancestor order.
The resulting tree has at most \(\bar L\le\bar L_A\) leaves.  Each snapped
boundary disagrees with the truth only in a slab of width \(O(r^{-1})\).
By (J) and the upper density bound in (Dens), the loss contributed by one
such slab is at most
\(\kappa_{\max}^2f_{\max}/r\), after absorbing the fixed quasi-uniformity
constant \(\gamma_{\mathrm{grid}}\) into \(f_{\max}\).  Since a tree with at most \(\bar L\) leaves has
at most \(\bar L-1\) boundaries,
\[
  \min_{\sigma\in\cT_{\bar L_A}(\cA_r)}R(\sigma)
  -
  R(\gamma_0)
  \le
  \frac{(\bar L-1)\kappa_{\max}^2f_{\max}}{r}.
\]
Subtracting this upper bound from the omitted-boundary lower bound (using
the halved \(c_{\mathrm{topol}}/2\) from the discard adjustment above) yields
\[
  S_r(\tau)
  \ge
  \frac{
    \frac12 c_{\mathrm{topol}}\kappa_{\min}^2f_{\min}q_{\min}C
    -
    (\bar L-1)\kappa_{\max}^2f_{\max}
  }{r}
  =
  \frac{c_\Delta}{r}.
\]
Condition~\eqref{eq:C-margin-choice} implies not only \(c_\Delta>0\), but
\[
  c_\Delta
  >
  \frac14c_{\mathrm{topol}}\kappa_{\min}^2f_{\min}q_{\min}C.
\]
This completes the proof.
\end{proof}

\begin{proof}[Proof of Proposition~\ref{prop:parsimony-grid}]
\emph{Proof strategy.}
We construct a feasible grid tree representing the atom-projected truth,
control all atom-mean noise by the single quadratic quantity \(V_r\), and
separate that noise from empirical occupancy by introducing the
empirical-design signal \(\widehat S_r\).  A local relative-occupancy argument
gives an \(r^{-1}\) lower bound for omitted boundaries.  The penalty then
excludes both omissions and unmatched extra splits.

For brevity write \(r=r_n\), \(N_r=N_{r_n}\), and
\(\gbar_a=\gbar_r(a)\).  Let \(\tau_r^\star\) be the tree obtained by
placing two consecutive grid splits around every true boundary lying inside
an atom and one split at every true boundary already lying on a cutpoint.
It has at most
\[
  \bar L+(\bar L-1)=2\bar L-1
\]
leaves and depth at most \(2d\), and it is constant on the level sets of
\(\gbar_r\).  Therefore
\[
  \tau_r^\star\in\cS^{\cA_r},
  \qquad
  \Pi^\nu_{\tau_r^\star}\gbar_r=\gbar_r.
\]
Under \(m_n=o(n/r_n)\), every transition leaf has expected occupancy
asymptotically larger than \(m_n\); the finitely many remaining leaves have
fixed positive mass by (Node).  Binomial concentration therefore makes
\(\tau_r^\star\) Stage-A feasible with probability tending to one.

For each atom \(a\in\cA_r\), let \(\bar Z_a\) be its empirical response mean
and define
\[
  D_a:=\bar Z_a-\gbar_a.
\]
When an atom is empty, set \(D_a=0\); its empirical weight is then zero.
Define
\[
  V_r:=\sum_{a\in\cA_r}\widehat\omega_aD_a^2.
\]
Conditional on a positive atom occupancy \(n_a\),
Assumption~\ref{ass:moments} gives
\[
  \E[D_a^2\mid n_a]\le\frac{\sigma_a^2}{n_a}
  \le\frac{\sigma^2}{n_a},
  \qquad
  \sigma^2:=\sup_a\sigma_a^2<\infty.
\]
Since \(\widehat\omega_a=n_a/n\) on the corresponding unnormalised risk
scale,
\[
  \E[\widehat\omega_aD_a^2\mid n_a]
  \le
  \frac{n_a}{n}\frac{\sigma^2}{n_a}
  =
  \frac{\sigma^2}{n}
\]
when \(n_a>0\), while the left-hand side is zero when \(n_a=0\).
Taking expectations and summing gives the exact inequality
\[
  \E[V_r]\le\frac{N_r\sigma^2}{n}.
\]
Markov's inequality therefore yields
\begin{equation}\label{eq:Vr-rate}
  V_r=O_p(N_r/n)=o_p(\lambda_n)=o_p(r^{-1}).
\end{equation}
No homoskedasticity assumption, union bound, or global lower bound on all
atom occupancies is used.

For \(\tau\in\cT_{\bar L_A}(\cA_r)\), let \(\cG_\tau\) be the linear space of
functions constant on the leaves of \(\tau\), and define the
empirical-design signal
\begin{equation}\label{eq:empirical-design-signal}
  \widehat S_r(\tau)
  :=
  \min_{g\in\cG_\tau}
  \sum_{a\in\cA_r}
  \widehat\omega_a\{\gbar_a-g(a)\}^2.
\end{equation}
This is the approximation error for the population atom means
\(\gbar_a\), but under the empirical atom weights.  It therefore incorporates
occupancy imbalance while excluding response-mean noise.

Let \(\widehat\Pi_\tau\) denote orthogonal projection under the empirical
inner product
\[
  \langle u,v\rangle_{\widehat\omega}
  :=
  \sum_{a\in\cA_r}\widehat\omega_a u(a)v(a).
\]
Then
\[
  \widehat S_r(\tau)
  =
  \|(I-\widehat\Pi_\tau)\gbar\|_{\widehat\omega}^2.
\]
The empirical ANOVA identity gives
\[
  R_n(\tau)-R_n(\tau^{\mathrm{sat}}_r)
  =
  \|(I-\widehat\Pi_\tau)(\gbar+D)\|_{\widehat\omega}^2.
\]
Because \(\tau_r^\star\) represents \(\gbar\),
\[
  \|(I-\widehat\Pi_{\tau_r^\star})\gbar\|_{\widehat\omega}=0
\]
and, since orthogonal projections are contractions,
\[
  R_n(\tau_r^\star)-R_n(\tau^{\mathrm{sat}}_r)
  =
  \|(I-\widehat\Pi_{\tau_r^\star})D\|_{\widehat\omega}^2
  \le V_r.
\]
The reverse triangle inequality and contraction of
\(I-\widehat\Pi_\tau\) give
\begin{align*}
  \|(I-\widehat\Pi_\tau)(\gbar+D)\|_{\widehat\omega}
  &\ge
  \|(I-\widehat\Pi_\tau)\gbar\|_{\widehat\omega}
  -
  \|(I-\widehat\Pi_\tau)D\|_{\widehat\omega}\\
  &\ge
  \sqrt{\widehat S_r(\tau)}-\sqrt{V_r}.
\end{align*}
Squaring the nonnegative lower bound when
\(\widehat S_r(\tau)\ge V_r\), and using the trivial lower bound zero
otherwise, yields uniformly
\[
  \|(I-\widehat\Pi_\tau)(\gbar+D)\|_{\widehat\omega}^2
  \ge
  \widehat S_r(\tau)
  -2\sqrt{\widehat S_r(\tau)V_r}.
\]
Consequently,
\begin{equation}\label{eq:correct-empirical-sandwich}
  R_n(\tau)-R_n(\tau_r^\star)
  \ge
  \widehat S_r(\tau)
  -2\sqrt{\widehat S_r(\tau)V_r}
  -V_r.
\end{equation}
Unlike the corresponding formula with the population-weighted
\(S_r(\tau)\), this identity contains no population-versus-empirical weight
mismatch: the deterministic signal, cross term, and noise norm all use the
same empirical inner product.

We next transfer the omitted-boundary margin to the empirical design.
For each of the finitely many true boundaries, use the two fixed local
half-patches constructed in Lemma~\ref{lem:topol-margin}.  Their population
masses are bounded above and below by fixed multiples of \(r^{-1}\).
The intersections of these patches with leaves of trees having at most
\(\bar L_A\) leaves form a fixed-VC-dimension family of axis-aligned sets.
The local relative Bernstein bound for this family gives
\[
  \sup_B
  \frac{\abs{\widehat\nu(B)-\nu(B)}}
       {\nu(B)}
  =
  o_p(1)
\]
uniformly over the relevant boundary-patch intersections having mass of
order \(r^{-1}\), because
\[
  \frac{r\log n}{n}\longrightarrow0,
\]
which follows from \(r^{d+1}=o(n)\): the latter gives \(r=o(n^{1/(d+1)})\),
so \(r\log n=o(n^{1/(d+1)}\log n)=o(n)\) for every fixed \(d\ge1\) (the
relative-deviation bound for a mass-\(\Theta(r^{-1})\), fixed-VC-dimension
family requires \(r\log n/n\to0\), not merely \(r\log r/n\to0\), since the
uniform bound is taken over \(n\)-dependent occupancy counts, not over the
\(r\) atoms themselves).  This is a local boundary-patch
occupancy statement, not global control of all \(N_r\) atom occupancies.
Since the number of true boundaries and active face cells is fixed, the
pairing calculation of Lemma~\ref{lem:topol-margin}, now with empirical
weights, implies that for some fixed \(c_\Delta'>0\),
\begin{equation}\label{eq:empirical-local-margin}
  \PP\left\{
    \inf_{\tau:\,\tau\text{ omits a true boundary}}
    \widehat S_r(\tau)
    \ge\frac{c_\Delta'}{r}
  \right\}
  \longrightarrow1.
\end{equation}
For example, after increasing \(r_0\), one may take
\[
  c_\Delta'
  =
  \frac14c_{\mathrm{topol}}\kappa_{\min}^2f_{\min}q_{\min}C.
\]

We now exclude the two possible topology errors.

\emph{Omitted boundaries.}
On the event in~\eqref{eq:empirical-local-margin}, let \(\tau\) omit a true
boundary.  By~\eqref{eq:Vr-rate},
\[
  \frac{V_r}{\widehat S_r(\tau)}
  \le
  \frac{rV_r}{c_\Delta'}
  =o_p(1)
\]
uniformly over all such \(\tau\).  Hence, with probability tending to one,
\(V_r\le\widehat S_r(\tau)/16\), and
\begin{align*}
  R_n(\tau)-R_n(\tau_r^\star)
  &\ge
  \widehat S_r(\tau)
  -2\sqrt{\widehat S_r(\tau)V_r}
  -V_r\\
  &\ge
  \widehat S_r(\tau)
  -\frac12\widehat S_r(\tau)
  -\frac1{16}\widehat S_r(\tau)\\
  &=
  \frac7{16}\widehat S_r(\tau)\\
  &\ge
  \frac{7c_\Delta'}{16r}.
\end{align*}
The largest possible penalty advantage of \(\tau\) relative to
\(\tau_r^\star\) is at most \(\lambda_n\bar L_A=o(r^{-1})\).  Therefore
\[
  R_n(\tau)+\lambda_n\card\tau
  >
  R_n(\tau_r^\star)
  +\lambda_n\card{\tau_r^\star}
\]
uniformly over omitted-boundary candidates, with probability tending to one.

\emph{Unmatched extra splits.}
Conditional on retaining at least one representative of every true
boundary, the full-support condition and (Local-Jump) imply the following
deterministic reduction.  Any collapsed split not associated with a true
boundary lies wholly inside a region on which \(\gbar_r\) is constant.
Contract that split and rebuild its descendants in the true ancestor order.
Repeating this operation produces a tree \(\tau^{-}\) with
\[
  \card{\tau^{-}}\le\card\tau-1,
  \qquad
  \Pi^\nu_{\tau^{-}}\gbar_r
  =
  \Pi^\nu_\tau\gbar_r.
\]
Transition pairs enclosing true boundaries are retained and count as one
collapsed boundary.  Thus the contraction removes only unmatched
topological splits and does not remove a transition split needed to represent
a projected jump.

Because \(\tau\) and \(\tau^{-}\) have the same population-mean fit, any
empirical improvement of \(\tau\) over \(\tau^{-}\) is caused solely by
\(D\).  Write \(\ell\) for the contracted leaf of \(\tau^-\) and
\(\ell_1,\ell_2\) for its two children in \(\tau\) (the argument iterates
over every contracted split); since \(\gbar_r\) is constant
\(\gbar_\ell\) on \(\ell\), \(\widehat\gamma_{\ell_k}-\widehat\gamma_\ell
=\bar D_{\ell_k}-\bar D_\ell\), where \(\bar D_S\) denotes the
\(\widehat\omega\)-weighted average of \(D_a\) over atoms \(a\subseteq S\).
By \((x-y)^2\le2x^2+2y^2\) and Jensen's inequality applied to the two
weighted averages,
\[
  \widehat\omega_{\ell_k}(\widehat\gamma_{\ell_k}-\widehat\gamma_\ell)^2
  \le
  2\sum_{a\subseteq\ell_k}\widehat\omega_aD_a^2
  +2\widehat\omega_{\ell_k}\,\overline{D_a^2}_\ell,
\]
and summing over \(k=1,2\) collapses the second term to
\(2\sum_{a\subseteq\ell}\widehat\omega_aD_a^2\) exactly, by definition of
the weighted average over \(\ell\).  Orthogonal-projection contraction and
Cauchy--Schwarz therefore give
\[
  R_n(\tau^{-})-R_n(\tau)
  =
  \sum_{k=1,2}\widehat\omega_{\ell_k}
    (\widehat\gamma_{\ell_k}-\widehat\gamma_\ell)^2
  \le
  4\sum_{a\subseteq\ell}\widehat\omega_aD_a^2
  \le 4V_r.
\]
Accordingly,
\begin{align*}
  &\bigl\{R_n(\tau)+\lambda_n\card\tau\bigr\}
  -
  \bigl\{R_n(\tau^{-})
          +\lambda_n\card{\tau^{-}}\bigr\}\\
  &\qquad\ge
  \lambda_n-4V_r
  =
  \lambda_n\{1-o_p(1)\}
  >0
\end{align*}
with probability tending to one, by \(V_r=o_p(\lambda_n)\).  Hence no
penalised minimiser contains an unmatched collapsed split.

Combining the two exclusions, a Stage-A minimiser omits no true boundary and
contains no unmatched collapsed split.  Its collapsed split set is therefore
in bijection with the true boundaries.  By the definition of omission, the
representative paired with each true boundary lies within \(C/r_n\) of it.
\end{proof}

\begin{proof}[Proof of Proposition~\ref{prop:boundary-recovery}]
\emph{Proof strategy.}
Apply Proposition~\ref{prop:parsimony-grid} along the joint sequence
\((n,r_n,\lambda_n,m_n)\); no fixed-\(r\), then \(n\to\infty\), interchange is
used.

All hypotheses of Proposition~\ref{prop:parsimony-grid} hold along the stated
sequence.  Its conclusion therefore gives the asserted bijection and
\(C/r_n\) localisation with probability tending to one.

To verify that this is the conclusion required by Stage~B, consider two
distinct true boundaries on the same coordinate.  By (Sep), their distance
is at least \(\varsigma\).  Their recovered representatives are each within
\(C/r_n\) of their true values, so the triangle inequality gives a recovered
separation of at least
\[
  \varsigma-\frac{2C}{r_n}
  \longrightarrow\varsigma.
\]
Thus, eventually, the bracket bounded by the nearest \emph{other}
same-coordinate recovered split has width bounded below by a fixed positive
constant, for example \(\varsigma/2\).  The Stage-B bracket does not shrink
at rate \(r_n^{-1}\); it remains of fixed asymptotic width while containing
the true boundary.  This is exactly the topology-recovery conclusion used by
the separately established Stage-B rate proposition.
\end{proof}

\begin{proof}[Proof of Proposition~\ref{prop:transition-leaves}]
A boundary at $t^\star$ interior to an atom $a$ makes $\gbar_r$ take on $a$
an intermediate value distinct from its values on the neighbouring atoms, by
clause (J) and the positive mass of both parts of $a$ under (Dens). By
Lemma~\ref{lem:grid-optimal} a grid-optimal tree must therefore separate $a$
from both neighbours, costing one leaf beyond the $L^{\mathrm{cont}}$ leaves
of the snapped partition. Summing over the at most $L^{\mathrm{cont}}-1$
boundaries gives the outer bounds; boundaries already at a cutpoint cost
nothing. Define the collapse map $\topolcoll$, merging any two splits on one
coordinate at consecutive cutpoints of $\cA_r$ within a node into the single
interval they span; $\topolcoll$ of a grid-optimal tree recovers the true
topology (Proposition~\ref{prop:boundary-recovery} makes this precise, with
the $C/r$ tolerance), and the interval it returns is exactly the enclosing
topological interval on which Stage~B searches.
\end{proof}

\refstepcounter{subsection}
\subsection*{S.3\quad Proofs for Section~\ref{sec:stageb}}%
\label{sec:supp-stageb}
\addcontentsline{toc}{subsection}{S.3\quad Proofs for
  Section~\protect\ref{sec:stageb}}

\begin{proof}[Proof of Lemma~\ref{lem:bracket-reduction}]
\emph{Proof strategy.} Part (a) is a strict-inequality/compactness argument
using only population quantities. Part (b) upgrades it to consistency of
the empirical minimiser via a uniform law of large numbers on the fixed
compact $[q_N^-,q_N^+]$. Part (c) then uses (b) together with
Proposition~\ref{prop:boundary-recovery}'s bracket-width guarantee to show
the actual (random-bracket) minimiser and the fixed-bracket minimiser
coincide with probability tending to one.

\emph{(a).} For every $c\in[q_N^-,q_N^+]$,
$R_N^X(c)-R_N^X(c_0)=\|\Pi^\nu_{\tau(c)}\gamma_0-\gamma_0\|_\nu^2\ge0$,
the population ANOVA identity, with equality if and only if $\tau(c)$ is
sufficient for $\gamma_0$ (i.e.\ leaf-constant for it). For $c\ne c_0$, the
slab $\mathcal I_N(c,c_0)$ has strictly positive $\nu$-mass by (Dens), and by
(Local-Jump) at least one of its cells carries a nonzero jump; hence
$\tau(c)$ is not sufficient and the inequality is strict for every
$c\in[q_N^-,q_N^+]\setminus\{c_0\}$. By (Dens), $R_N^X$ is continuous in
$c$ on $[q_N^-,q_N^+]$. For every fixed $\eta\in(0,q_N^+-q_N^-]$, a
strictly positive continuous function on the
compact set $\{c\in[q_N^-,q_N^+]:\abs{c-c_0}\ge\eta\}$ attains a
strictly positive infimum, giving (a).

\emph{(b).} Let $\widehat c_N^{\mathrm{full}}$ denote the minimiser of the
empirical criterion over the FIXED, deterministic set $[q_N^-,q_N^+]$ ---
an auxiliary statistic, distinct from $\widehat c_N$ (the actual Stage-B
output, minimising over the random $\widehat B_n$), introduced only for
this proof. The empirical criterion depends on $c$ only through the
finitely many leaf or prism-cell counts and response sums, each of which
is a partial sum indexed by the single indicator $\mathbf 1\{X_q\le c\}$,
monotone in $c$. Uniform convergence of each such partial sum to its
population counterpart over the fixed compact $[q_N^-,q_N^+]$ follows from
the classical Glivenko--Cantelli theorem together with
Assumption~\ref{ass:moments}'s bound on the response second moment.
Combined with (a), the standard argmin-consistency argument for
M-estimators on a fixed compact parameter space
(\citealt{vandervaart1998}) gives
\[
  \widehat c_N^{\mathrm{full}}\to_pc_0 .
\]

\emph{(c).} By Proposition~\ref{prop:boundary-recovery}, the recovered
same-coordinate separation around $c_0$ is at least $\varsigma-2C/r_n$,
which exceeds $\varsigma/2$ for $n$ large; consequently
\[
  \PP\bigl(B^-\subseteq\widehat B_n\bigr)\longrightarrow1,
\]
while $\widehat B_n\subseteq[q_N^-,q_N^+]$ holds always, since Stage B
never searches outside the node's own region. By (b),
$\PP(\widehat c_N^{\mathrm{full}}\in B^-)\to1$. On the intersection of
these two probability-tending-to-one events:
$\widehat c_N^{\mathrm{full}}$ is the global minimiser of the criterion
over $[q_N^-,q_N^+]$, and $B^-\subseteq\widehat B_n\subseteq
[q_N^-,q_N^+]$, with $\widehat c_N^{\mathrm{full}}$ landing in the smallest
of these three nested sets, $B^-$. A global minimiser over a set that
lands inside a subset of that set is automatically the minimiser over
every intermediate subset containing it; applying this twice,
$\widehat c_N^{\mathrm{full}}$ is simultaneously the minimiser over
$\widehat B_n$ and over $B^-$. The former identifies
$\widehat c_N^{\mathrm{full}}=\widehat c_N$ (both minimise the same
criterion over $\widehat B_n$); the latter identifies
$\widehat c_N^{\mathrm{full}}=\widetilde c_N$ (both minimise the same
criterion over $B^-$). Hence $\widehat c_N=\widetilde c_N$ on this
intersection, and $\PP(\widehat c_N=\widetilde c_N)\to1$. Combined with
(b), this also gives $\widehat c_N\to_pc_0$, proving the consistency
statement of part (b) for the actual Stage-B output.
\end{proof}

\begin{proof}[Proof of Lemma~\ref{lem:descendant-stability}]
Each descendant threshold displacement of size \(O(1/r_n)\) alters the
\(\nu\)-mass of the affected leaf regions by \(O(1/r_n)\), by the upper
density bound of (Dens) and the bounded coordinate diameters, so (a) follows
from the boundedness of the profiled leaf means
(Assumption~\ref{ass:design}, Assumption~\ref{ass:construct}(c)). Part (b)
follows from (a) and the
strengthened Lemma~\ref{lem:bracket-reduction}(a), for \(r_n\) large: the
perturbed infimum is within \(O(1/r_n)\) of the true one, hence exceeds
\(\delta_N(\eta)/2\) once \(1/r_n<\delta_N(\eta)/2\). Part (c) follows
because each perturbed adjacent-leaf pair differs from a true adjacent pair
only through leaf regions of \(\nu\)-mass \(O(1/r_n)\), whose profiled means
therefore differ by \(O(1/r_n)\) given the leaf-mass floor of (Node); with
\(r_n\to\infty\), (Local-Jump) transfers to the perturbed pairs with
\(\kappa_{\min}\) replaced by \(\kappa_{\min}/2\) for \(n\) large enough.
\end{proof}

\begin{proof}[Proof of Lemma~\ref{lem:prism-profile}]
\emph{Proof strategy.}
We first identify the observations whose assignment changes.  We then refine
the moving slab into finitely many cells, one for each possible source--target
leaf pair.  A weighted ANOVA identity gives an exact first-order
decomposition of the excess risk.  Finally, the local empirical criterion is
represented as a finite sum of marked change-point processes with positive
aggregate drift.

\emph{Step 1: identify the moving slab.}
Suppose first that \(c>c_0\).  If \(x_q\le c_0\), then both thresholds send
\(x\) to the left child of \(N\).  If \(x_q>c\), both send \(x\) to the right
child.  If \(c_0<x_q\le c\), the true threshold sends \(x\) right whereas the
candidate sends it left.  Therefore,
\[
  \bigl\{x\in R_N:
    \text{the side of \(N\) differs under \(c\) and \(c_0\)}
  \bigr\}
  =
  R_N\cap\{c_0<x_q\le c\}.
\]
The case \(c<c_0\) is identical with left and right reversed, proving (i).

\emph{Step 2: form the prism cells.}
Within the bracket, no candidate threshold crosses another split on
coordinate \(q\), by Assumption~\ref{ass:solver-twostage}(b) (which defines
the bracket as bounded by the nearest other same-coordinate split) and (Sep).
Hence every descendant test on coordinate \(q\) has the same outcome
throughout the slab and may be absorbed into a fixed path label.  The
remaining descendant tests are on coordinates other than \(q\), and their
outcomes are unaffected by moving \(c\).  Thus the other-coordinate values of
an observation determine exactly one leaf below the true side of \(N\) and
exactly one leaf below the candidate side.

Intersecting the corresponding side-specific leaf regions partitions the
slab into disjoint cells indexed by realised source--target pairs.  There are
at most \(L_N^-L_N^+\) such pairs.  Since
\(L_N^-+L_N^+\le\bar L\), the arithmetic--geometric mean inequality gives
\[
  L_N^-L_N^+
  \le
  \left(\frac{L_N^-+L_N^+}{2}\right)^2
  \le \frac{\bar L^2}{4}.
\]
This proves (ii).  By (Dens), each active cell has mass
\[
  \omega_{N,r}(c)
  =
  f_{N,r}^{\operatorname{sgn}(c-c_0)}(c_0)
  \abs{c-c_0}
  +o\!\left(\abs{c-c_0}\right),
\]
and (Local-Jump) gives
\(\kappa_{\min}\le\abs{\Delta_{N,r}}\le\kappa_{\max}\).
This proves (iii).

\emph{Step 3: derive the population excess risk.}
Consider one destination leaf.  Let its correctly assigned core have mass
\(a>0\) and population mean \(\mu_0\).  Suppose that the moving prism cells
entering this leaf have masses \(b_1,\ldots,b_s\) and means
\(\mu_1,\ldots,\mu_s\).  Put
\[
  B=\sum_{r=1}^s b_r,
  \qquad
  \Delta_r=\mu_r-\mu_0.
\]
The candidate leaf mean is
\[
  \bar\mu
  =
  \frac{a\mu_0+\sum_{r=1}^s b_r\mu_r}{a+B}
  =
  \mu_0+\frac{\sum_{r=1}^s b_r\Delta_r}{a+B}.
\]
The weighted ANOVA identity states that the excess squared error from pooling
groups with masses \(w_g\) and means \(\mu_g\) equals
\(\sum_g w_g(\mu_g-\bar\mu)^2\).  Applying it here and expanding gives
\begin{align*}
  &a(\mu_0-\bar\mu)^2
    +\sum_{r=1}^s b_r(\mu_r-\bar\mu)^2
  \\
  &\quad=
    \sum_{r=1}^s b_r\Delta_r^2
    -
    \frac{\left(\sum_{r=1}^s b_r\Delta_r\right)^2}{a+B}.
\end{align*}
Indeed,
\begin{align*}
  \sum_{g}w_g(\mu_g-\bar\mu)^2
  &=
  \sum_g w_g\mu_g^2-(a+B)\bar\mu^2
  \\
  &=
  a\mu_0^2+\sum_r b_r(\mu_0+\Delta_r)^2
  -
  \frac{\left\{(a+B)\mu_0+\sum_r b_r\Delta_r\right\}^2}{a+B}
  \\
  &=
  \sum_r b_r\Delta_r^2
  -
  \frac{\left(\sum_r b_r\Delta_r\right)^2}{a+B}.
\end{align*}
By (Node), \(a\ge q_{\min}-B\), and hence \(a\ge q_{\min}/2\) for all
sufficiently small \(\abs{c-c_0}\).  Moreover,
\[
  0
  \le
  \frac{\left(\sum_r b_r\Delta_r\right)^2}{a+B}
  \le
  \frac{\kappa_{\max}^2 B^2}{a+B}
  \le
  \frac{2\kappa_{\max}^2}{q_{\min}}B^2.
\]
Summing over the finitely many destination leaves, and using
\(B=O(\abs{c-c_0})\), yields
\[
  R_N^{\mathrm{prof}}(c)-R_N^{\mathrm{prof}}(c_0)
  =
  \sum_r\omega_{N,r}(c)\Delta_{N,r}^2
  +O\!\left(\abs{c-c_0}^2\right),
\]
which is \eqref{eq:prism-risk}.  Equivalently,
\[
  R_N^{\mathrm{prof}}(c)-R_N^{\mathrm{prof}}(c_0)
  =
  \left\{\sum_r\omega_{N,r}(c)\Delta_{N,r}^2\right\}
  \{1+o(1)\}.
\]

By (Dens),
\[
  \sum_r\omega_{N,r}(c)
  \ge \frac{p_{\min}}{2}\abs{c-c_0}
\]
for all sufficiently small \(\abs{c-c_0}\).  Therefore,
\[
  \sum_r\omega_{N,r}(c)\Delta_{N,r}^2
  \ge
  \frac{\kappa_{\min}^2p_{\min}}{2}\abs{c-c_0}.
\]
Shrinking the neighbourhood once more so that the quadratic remainder is no
larger than half this leading term gives
\[
  R_N^{\mathrm{prof}}(c)-R_N^{\mathrm{prof}}(c_0)
  \ge
  \frac{\kappa_{\min}^2p_{\min}}{4}\abs{c-c_0},
\]
proving (iv).

\emph{Step 4: identify the local empirical process.}
Let \(c=c_0+u/n\), with \(u\) in a fixed compact set.  The expected number of
observations in each active prism cell is
\[
  n\omega_{N,r}(c)
  =
  \abs{u}\,
  f_{N,r}^{\operatorname{sgn}(u)}(c_0)+o(1),
\]
and is therefore of constant order.  If an observation moves from a leaf
with mean \(\mu_s\) to one with mean \(\mu_d\), write
\(Z=\mu_s+\varepsilon\), where
\(\E(\varepsilon\mid X)=0\).  Its unnormalised SSE increment is
\begin{align*}
  (Z-\mu_d)^2-(Z-\mu_s)^2
  &=
  (\varepsilon+\mu_s-\mu_d)^2-\varepsilon^2
  \\
  &=
  2\varepsilon(\mu_s-\mu_d)
  +(\mu_s-\mu_d)^2.
\end{align*}
Conditional on its prism cell, this increment has mean
\[
  \E\!\left[
    (Z-\mu_d)^2-(Z-\mu_s)^2
    \mid X
  \right]
  =
  (\mu_s-\mu_d)^2
  =
  \Delta_{N,r}^2
  \ge\kappa_{\min}^2,
\]
and finite variance by the standing conditional second-moment bound.

Consequently, the local unnormalised criterion is a finite sum of marked
rare-event counting processes, all indexed by the same scalar \(u\), with
aggregate one-sided mean drift at least
\(\kappa_{\min}^2p_{\min}\abs u/2\).  The number of component processes is
bounded by \(\bar L^2/4\), independently of \(n\).  Joint convergence follows
from the standard finite-dimensional marked point-process argument: the
disjoint prism-cell counts converge jointly to Poisson counts with their
one-sided limiting intensities, while the associated marks have finite second
moments.  Summation is a continuous linear map on this finite vector of
processes.  Thus the sum has the same \(n\)-localisation as a single
change-point process.

This is the random-design extension of the single-break mechanism in
\citet{chan1993} and \citet{seijosen2011}.  The simultaneous-break
least-squares machinery of \citet{baiperron1998} provides the corresponding
finite-dimensional profiling precedent, while \citet{yaoau1989} and
\citet{boysen2009} establish the same \(n^{-1}\) location rate when finitely
many step heights are estimated simultaneously.  Those references do not by
themselves identify the present leaf-pair sum; the prism decomposition and
joint marked-process argument above supply that missing reduction.

Step 4 is a purely LOCAL statement, valid for \(c=c_0+u/n\) with \(u\) in a
fixed compact set: on its own it bounds the criterion's behaviour near
\(c_0\) but does not rule out the empirical minimiser settling elsewhere in
\(B^-\).  That is supplied separately, and globally, by
Lemma~\ref{lem:bracket-reduction}: applying it with \(R_N^X=R_N^{\mathrm
{prof}}\) gives \(\widetilde c_N\to_pc_0\) (part (b), via the same
Glivenko--Cantelli argument used there, now applied to the profiled
criterion's finitely many leaf/cell sufficient statistics). Consistency and
the local drift bound~\eqref{eq:prism-drift} together license the standard
non-standard-rate argmax theorem for change-point processes
(\citealt{seootsu2018} give a general version covering discontinuous,
possibly localised criteria of exactly this kind): consistency confines the
search to shrinking neighbourhoods of \(c_0\), on which Step 4's local
process governs the limit.

Because the limiting two-sided process has positive linear drift away from
zero, its set of minimisers is tight.  The standard change-point argmin
argument therefore yields
\[
  \lim_{U\to\infty}\limsup_{n\to\infty}
  \PP\!\left(n\abs{\widetilde c_N-c_0}>U\right)=0,
\]
which is equivalent to
\(n(\widetilde c_N-c_0)=O_p(1)\), where \(\widetilde c_N\) is the minimiser
of the profiled criterion over the FIXED interval \(B^-\) of
Lemma~\ref{lem:bracket-reduction}.  By that lemma's part (c),
\(\PP(\widehat c_N=\widetilde c_N)\to1\), so the rate transfers to the
actual Stage-B output: \(n(\widehat c_N-c_0)=O_p(1)\).  This proves (v).
\end{proof}

\begin{proof}[Proof of Lemma~\ref{lem:ancestor-localisation}]
\emph{Proof strategy.}
We localise every ancestor error to a deterministic \(M/n\) collar.  A joint
density union bound shows that, with probability tending to one, no
ancestor-contamination observation also lies in the level-\(\ell\)
\(U/n\)-window.  On that event, contaminated observations have assignments
that are constant in the local candidate threshold and perturb each profiled
cell criterion by \(o_p(1)\) on the unnormalised local-SSE scale.

\emph{Step 1: deterministic ancestor collars.}
Fix \(\epsilon>0\), and let
\(\operatorname{anc}(\ell)\) be the finite set of ancestors of the
level-\(\ell\) node.  By the induction hypothesis, there is \(M<\infty\) such
that, for every \(a\in\operatorname{anc}(\ell)\),
\[
  \PP\!\left(
    n\abs{\widehat c_a-c_{0,a}}>M
  \right)
  <
  \frac{\epsilon}
       {3\card{\operatorname{anc}(\ell)}}.
\]
Hence, by the union bound,
\[
  \PP\!\left(
    \max_{a\in\operatorname{anc}(\ell)}
    n\abs{\widehat c_a-c_{0,a}}>M
  \right)
  <\frac{\epsilon}{3}.
\]
On the complementary event, every observation whose membership differs
between true and refined ancestor conditioning lies in
\[
  \mathcal C_{n,M}
  :=
  \bigcup_{a\in\operatorname{anc}(\ell)}
  \left\{
    x:\abs{x_{q_a}-c_{0,a}}\le\frac{M}{n}
  \right\}.
\]

\emph{Step 2: exclude overlap with the descendant localisation window.}
Fix \(U<\infty\).  If \(q_a\ne q_\ell\), clause (Joint) and a bare union bound
give
\begin{align*}
  &\PP\!\left(
    \exists i:
    \abs{X_{i,q_a}-c_{0,a}}\le\frac{M}{n},
    \ 
    \abs{X_{i,q_\ell}-c_{0,\ell}}\le\frac{U}{n}
  \right)
  \\
  &\quad\le
  n\,
  \PP\!\left(
    \abs{X_{q_a}-c_{0,a}}\le\frac{M}{n},
    \ 
    \abs{X_{q_\ell}-c_{0,\ell}}\le\frac{U}{n}
  \right)
  \\
  &\quad\le
  n\left(\frac{2M}{n}\right)
       \left(\frac{2U}{n}\right)\bar f_2
  =
  \frac{4\bar f_2MU}{n}
  \longrightarrow0.
\end{align*}
A further finite union bound handles all ancestors.  If
\(q_a=q_\ell\), (Sep) gives
\(\abs{c_{0,a}-c_{0,\ell}}\ge\varsigma\); hence the two collars are disjoint
for all \(n>2(M+U)/\varsigma\), and the corresponding bad event is empty.

It follows that, with probability tending to one, every
ancestor-contamination observation lies outside
\[
  \left\{
    x:\abs{x_{q_\ell}-c_{0,\ell}}\le U/n
  \right\}.
\]
The number of observations in \(\mathcal C_{n,M}\) is \(O_p(1)\), because its
expectation is bounded uniformly in \(n\) by the marginal upper-density bound
and the fixed number of ancestors.

\emph{Step 3: bound the criterion perturbation.}
On the preceding good event, each contaminated observation enters or leaves
a specific prism cell in the level-\(\ell\) subtree, but its cell membership
does not change as \(c\) ranges over
\(c_{0,\ell}+[-U,U]/n\).  Thus it changes the count and response sum of each
affected cell by \(O_p(1)\).  By Assumption~\ref{ass:solver-twostage}(g),
every leaf of the refined partition contains at least \(\rho n\)
observations with probability tending to one, so its fitted mean changes by
\[
  O_p(1)/\Theta_p(n)=O_p(n^{-1}).
\]
The level-\(\ell\) local contrast itself involves only \(O_p(1)\) observations
whose \(q_\ell\)-coordinate lies in the \(U/n\)-window.  Differentiating a
single squared residual with respect to a leaf mean gives
\[
  (Z-\mu-\varrho)^2-(Z-\mu)^2
  =
  -2\varrho(Z-\mu)+\varrho^2.
\]
With \(\varrho=O_p(n^{-1})\), a tight number of local observations, and the
standing
conditional second-moment bound, summing this display gives \(O_p(n^{-1})\).
The corresponding change in the profiled least-squares gain is also
\(O_p(n^{-1})\), by the quadratic identity used in
Remark~\ref{rem:profiling-correction}.  Therefore,
\begin{align}
\label{eq:ancestor-op}
  \sup_{\abs u\le U}
  \Big|
  &\{S^{\mathrm{act}}_{n,\ell}(c_{0,\ell}+u/n)
     -S^{\mathrm{act}}_{n,\ell}(c_{0,\ell})\}
  \nonumber\\[-2pt]
  &-
  \{S^{\mathrm{or}}_{n,\ell}(c_{0,\ell}+u/n)
     -S^{\mathrm{or}}_{n,\ell}(c_{0,\ell})\}
  \Big|
  =
  O_p(n^{-1})
  =
  o_p(1),
\end{align}
where \(S^{\mathrm{act}}_{n,\ell}\) and \(S^{\mathrm{or}}_{n,\ell}\) are the
unnormalised actual-row-set and oracle-row-set subtree SSEs.

The oracle local process has a tight minimiser by
Lemma~\ref{lem:prism-profile}. Step~3 controls the actual--oracle
discrepancy only on a fixed compact $u$-window; tightness of the actual
process additionally requires that its criterion cannot be minimised far
from $c_{0,\ell}$, which is Step~4 below.

\emph{Step 4: exclude escape from the compact window.} Only observations in
the ancestor collar $\mathcal C_{n,M}$ can contribute to the actual--oracle
discrepancy outside the $U/n$-window used in Step~3; write
\[
  \Xi_n(u)
  :=
  \{S^{\mathrm{act}}_{n,\ell}(c_{0,\ell}+u/n)-S^{\mathrm{act}}_{n,\ell}(c_{0,\ell})\}
  -
  \{S^{\mathrm{or}}_{n,\ell}(c_{0,\ell}+u/n)-S^{\mathrm{or}}_{n,\ell}(c_{0,\ell})\}
\]
for $\abs u\le\varsigma n/4$.

\emph{Step 4a: the discrepancy is uniformly $O_p(1)$, with no $u$-dependence.}
As in Step~3, only two mechanisms contribute: observations whose
$q_\ell$-side membership changes between $c_{0,\ell}$ and $c_{0,\ell}+u/n$,
and the indirect effect through the profiled leaf means; the indirect
effect is $O_p(n^{-1})$ uniformly in $u$ by Step~3's own argument (the
$O_p(1)$ collar count and the $\Theta_p(n)$ leaf denominators of
Assumption~\ref{ass:solver-twostage}(g) do not depend on $u$), hence
$o_p(1)$ uniformly and negligible below. For the direct effect, every
contributing observation lies in $\mathcal C_{n,M}$ \emph{regardless of
$u$} --- the collar is defined by the ancestor coordinate alone, and only
an observation's \emph{side} assignment, never its collar membership,
depends on $u$. Write
\[
  K_n
  :=
  \sum_{i:X_i\in\mathcal C_{n,M}}
  \{Z_i-\gamma_0(X_i)\}^2
\]
for the aggregate squared residual over the collar. By the standing
conditional second-moment bound and $\PP(X\in\mathcal C_{n,M})=O(M/n)$
(Step~1),
\[
  \E[K_n]
  =
  n\,\PP(X\in\mathcal C_{n,M})\,
  \sup_x\E\bigl[\{Z-\gamma_0(X)\}^2\mid X=x\bigr]
  =
  O(M)=O(1),
\]
so $K_n=O_p(1)$ by Markov's inequality. The direct effect on $\Xi_n(u)$, at
any $u$, is a difference of squared-residual sums restricted to these same
collar observations, and is therefore bounded by a fixed multiple of $K_n$
\emph{uniformly in $u$}: no observation outside $\mathcal C_{n,M}$
contributes at any $u$, so there is nothing left to grow as $\abs u$ grows.
Hence
\[
  \sup_{\abs u\le\varsigma n/4}\abs{\Xi_n(u)}=O_p(1).
\]
This uniform bound needs no peeling and no $u$-indexed probability
calculation: once every contributing observation is confined, independently
of $u$, to the $O(1/n)$-mass collar $\mathcal C_{n,M}$, the range over which
$u$ varies cannot affect the bound.

\emph{Step 4b: the oracle criterion excludes escape.} The argmax-consistency
argument that establishes tightness for a profiled change-point criterion of
exactly this form --- via \citet{chan1993,seijosen2011} for Tier~1's
leaf-parent stump criterion, or via \citet{seootsu2018} for the profiled
criterion of part~(v) above --- proceeds by first showing that outside any
fixed neighbourhood of the true threshold, the criterion exceeds its
minimising value by an amount tending to infinity with the neighbourhood's
radius, uniformly, and only then concluding that the argmin is tight; the
tightness conclusion is a corollary of this escape-exclusion step, not an
independent fact. Lemma~\ref{lem:ancestor-localisation}'s own hypothesis
grants exactly this apparatus for $S^{\mathrm{or}}_{n,\ell}$: the
conclusions of Lemma~\ref{lem:prism-profile} (or its Tier-1 analogue) are
assumed to hold at level $\ell$ under true ancestor conditioning, which is
precisely $S^{\mathrm{or}}_{n,\ell}$. Applying the same escape-exclusion
step to $S^{\mathrm{or}}_{n,\ell}$ therefore gives: for every $K<\infty$,
there exists $U_K<\infty$ such that
\[
  \PP\!\left(
    \inf_{U_K\le\abs u\le\varsigma n/4}
    \bigl\{S^{\mathrm{or}}_{n,\ell}(c_{0,\ell}+u/n)
      -S^{\mathrm{or}}_{n,\ell}(c_{0,\ell})\bigr\}
    >K
  \right)
  \longrightarrow1.
\]

\emph{Step 4c: combine.} By Step~4a, choose $K_0<\infty$ with
$\PP\bigl(\sup_{\abs u\le\varsigma n/4}\abs{\Xi_n(u)}>K_0\bigr)<\epsilon/3$
for $n$ large. By Step~4b with $K=K_0$, choose $U\ge U_{K_0}$ with
\[
  \PP\!\left(
    \inf_{U\le\abs u\le\varsigma n/4}
    \bigl\{S^{\mathrm{or}}_{n,\ell}(c_{0,\ell}+u/n)
      -S^{\mathrm{or}}_{n,\ell}(c_{0,\ell})\bigr\}
    \le K_0
  \right)
  <\epsilon/3
\]
for $n$ large. On the intersection of these two probability-tending-to-one
events, for every $u$ with $U\le\abs u\le\varsigma n/4$,
\[
  S^{\mathrm{act}}_{n,\ell}(c_{0,\ell}+u/n)-S^{\mathrm{act}}_{n,\ell}(c_{0,\ell})
  \ \ge\
  \bigl\{S^{\mathrm{or}}_{n,\ell}(c_{0,\ell}+u/n)
    -S^{\mathrm{or}}_{n,\ell}(c_{0,\ell})\bigr\}
  -\abs{\Xi_n(u)}
  \ >\ K_0-K_0\ =\ 0,
\]
so the actual criterion strictly exceeds its value at $c_{0,\ell}$
throughout this range; since $\widehat c_\ell$ globally minimises
$S^{\mathrm{act}}_{n,\ell}$ over $B^-$ and $c_{0,\ell}\in B^-$ is a
competitor achieving increment exactly zero, $\widehat c_\ell$ cannot lie at
any such $u$, on an event of probability at least $1-2\epsilon/3$. The true
minimiser also cannot lie in $\{c:\varsigma n/4\le n\abs{c-c_{0,\ell}}\}$,
since $B^-$ already excludes it by construction. Together with Step~3's
control on $\abs u\le U$, this gives
$\PP(n\abs{\widehat c_\ell-c_{0,\ell}}>U)<\epsilon$ for $n$ large, completing
the third $\epsilon/3$ budgeted in Step~1. Hence
\[
  n(\widehat c_\ell-c_{0,\ell})=O_p(1).
\]
\end{proof}

\begin{proof}[Proof of Proposition~\ref{prop:twostage-superconsistency}]
\emph{Proof strategy.}
Apply the two-tier argument at the root and then proceed by induction over
tree levels, using Lemma~\ref{lem:ancestor-localisation} at every node with
ancestors.

The root rate follows from the classical stump argument when both children
are terminal and from Lemma~\ref{lem:prism-profile} otherwise.  If every
boundary above level \(\ell\) has the asserted rate,
Lemma~\ref{lem:ancestor-localisation} establishes it for every boundary at
level \(\ell\).  Since the depth is at most the fixed number \(d\), induction
covers the whole tree.

For each fixed \(k\) and every \(\epsilon>0\), choose \(M_k<\infty\) such that
\[
  \limsup_{n\to\infty}
  \PP\!\left(
    n\abs{\widehat c_k-c_{0,k}}>M_k
  \right)
  <
  \frac{\epsilon}{K}.
\]
With \(M=\max_{k\le K}M_k\), the finite union bound gives
\[
  \limsup_{n\to\infty}
  \PP\!\left(
    \max_{1\le k\le K}
    n\abs{\widehat c_k-c_{0,k}}>M
  \right)
  \le
  \sum_{k=1}^K
  \limsup_{n\to\infty}
  \PP\!\left(
    n\abs{\widehat c_k-c_{0,k}}>M
  \right)
  <\epsilon.
\]
Hence the maximum is \(O_p(1)\).  Since
\(K\le\bar L-1<\infty\), no growing-dimensional union bound is required.
\end{proof}

% ===========================================================================
\begin{thebibliography}{9}
% [TODO: port only the bibitems theory.tex \S6.1--6.3 actually cites --
% chan1993, seijosen2011, baiperron1998, yaoau1989, boysen2009, seootsu2018,
% vandervaart1998, xu2026 -- plus banerjee2007 for the related-work contrast
% in the introduction above. Do not port the causal-inference bibliography
% (chernozhukov2018, robins1995, hahn1998, van der Laan/Robins, etc.) --
% none of it is cited by the ported content.]
%
% Done for the ported content (2026-09-01). Every entry below is cited by
% Sections~\ref{sec:setup}--\ref{sec:stageb}; banerjee2007 is cited by the
% related-work paragraph in \S\ref{sec:setup-twostage} as well as by the
% introduction. xu2026 is NOT listed here: it is cited only by the
% introduction's own [TODO] outline, not by any ported content, and is left
% for whoever drafts \S\ref{sec:intro}.
\bibitem[Bai \& Perron(1998)]{baiperron1998} Bai, J. and Perron, P. (1998).
  Estimating and testing linear models with multiple structural changes.
  \emph{Econometrica} 66, 47--78.
\bibitem[Banerjee \& McKeague(2007)]{banerjee2007} Banerjee, M. and McKeague,
  I. W. (2007). Confidence sets for split points in decision trees.
  \emph{The Annals of Statistics} 35, 543--574.
\bibitem[Boysen et al.(2009)]{boysen2009} Boysen, L., Kempe, A., Liebscher,
  V., Munk, A. and Wittich, O. (2009). Consistencies and rates of convergence
  of jump-penalized least squares estimators. \emph{The Annals of Statistics}
  37, 157--183.
\bibitem[Chan(1993)]{chan1993} Chan, K. S. (1993). Consistency and limiting
  distribution of the least squares estimator of a threshold autoregressive
  model. \emph{The Annals of Statistics} 21, 520--533.
\bibitem[Seijo \& Sen(2011)]{seijosen2011} Seijo, E. and Sen, B. (2011).
  Change-point in stochastic design regression and the bootstrap.
  \emph{The Annals of Statistics} 39, 1580--1607.
\bibitem[Seo \& Otsu(2018)]{seootsu2018} Seo, M. H. and Otsu, T. (2018).
  Local M-estimation with discontinuous criterion for dependent and limited
  observations. \emph{The Annals of Statistics} 46, 344--369.
\bibitem[van der Vaart(1998)]{vandervaart1998} van der Vaart, A. W. (1998).
  \emph{Asymptotic Statistics}. Cambridge University Press.
\bibitem[Yao \& Au(1989)]{yaoau1989} Yao, Y.-C. and Au, S. T. (1989).
  Least-squares estimation of a step function. \emph{Sankhy\=a: The Indian
  Journal of Statistics, Series A} 51, 370--381.
\end{thebibliography}

\end{document}
